\documentclass[11pt]{amsart}
\usepackage[margin=1.1in]{geometry}
\usepackage{amsmath,amssymb,amsthm,mathtools}
\usepackage{booktabs}
\usepackage{microtype}
\usepackage{xcolor}
\usepackage[colorlinks=true,linkcolor=blue!50!black,citecolor=blue!50!black,urlcolor=blue!50!black]{hyperref}

\newtheorem{theorem}{Theorem}[section]
\newtheorem{proposition}[theorem]{Proposition}
\newtheorem{lemma}[theorem]{Lemma}
\newtheorem{corollary}[theorem]{Corollary}
\newtheorem{conjecture}[theorem]{Conjecture}
\theoremstyle{definition}
\newtheorem{definition}[theorem]{Definition}
\newtheorem{question}[theorem]{Question}
\theoremstyle{remark}
\newtheorem{remark}[theorem]{Remark}
\newtheorem{example}[theorem]{Example}

\newcommand{\Z}{\mathbb{Z}}
\newcommand{\F}{\mathbb{F}}
\newcommand{\A}{\mathcal{A}}
\newcommand{\Aut}{\operatorname{Aut}}
\newcommand{\End}{\operatorname{End}}
\newcommand{\Sym}{\operatorname{Sym}}
\newcommand{\Occ}{\operatorname{Occ}}
\newcommand{\Cent}{\mathcal{C}}
\newcommand{\sh}{\sigma}
\newcommand{\Lam}{\Lambda}

\title[The 88 second-order elementary cellular automata]{The 88 second-order elementary cellular automata: a classification by automorphism group}
\author{Michael Dayah}
\email{michael@dayah.com}

\begin{document}

\begin{abstract}
For an elementary rule $f$, the second-order automaton $F_f$ on $(\{0,1\}^2)^{\Z}$, given by $y_{t+1}=y_{t-1}\oplus f(y_t)$, is reversible, and the automorphism group of its spacetime subshift is the centralizer of $F_f$ in the automorphism group of the full $4$-shift. Up to reflection and the cellwise complement there are $88$ such automata, and we assign a status to each class. In $43$ classes the group contains every finite group; in $11$ it is $\langle\sh,F_f\rangle\cong\Z^2$; in $4$ it is $\langle\sh,F_f\rangle\times\Z/2$; in $2$ affine classes it is $\Z^2$ or $\Z^3$, strictly larger than $\langle\sh,F_f\rangle$; and $28$ classes are open. These entries rest on results of four companion papers: wall constructions, rigidity of the permutive rules and of the rules that read only two neighbours, edge moves for six classes, the rigidity theorem of Kitchens and Schmidt for the affine classes, and the groups of the classes of $94$ and $95$. We assemble them and add the following. An exhaustive SAT enumeration, repeated with several solvers, shows that in each open class every automorphism of radius at most $3$ lies in $\langle\sh,F_f\rangle$, and that for the rules $110$ and $30$ the block maps of radius at most $4$ that commute with $F_f$ are the maps $\sh^aF_f^b$ and two constants. The automaton $110$R has no wall of the known kinds, of any length. In at least $12$ classes, among them $10$ open ones, no configuration has a stuck edge, so the edge-move argument cannot start; computations, in part sampled, indicate where it stops in the other open classes.
\end{abstract}


\maketitle

\section{Introduction}\label{sec:intro}

Let $\A$ be a finite alphabet and $\sh$ the left shift on $\A^{\Z}$. By the Curtis--Hedlund--Lyndon theorem \cite{Hedlund1969} the endomorphisms of the full shift are the block maps, and its automorphisms the invertible block maps. The group $\Aut(\A^{\Z})$ is enormous: it contains every finite group and free groups \cite{BLR1988}, and every countable, locally finite, residually finite group \cite{KimRoush1990}, and its center consists of the powers of $\sh$ \cite{Ryan1972}. A reversible cellular automaton $F\in\Aut(\A^{\Z})$ determines a two-dimensional subshift, its \emph{spacetime} $X_F\subset\A^{\Z^2}$, whose rows are the configurations $F^n(c)$, $n\in\Z$ (see Cyr, Franks and Kra \cite{CFK2019} for the spacetime of a general endomorphism). The automorphism group of $X_F$ is the centralizer of $F$ in $\Aut(\A^{\Z})$ (Lemma~\ref{lem:spacetime}), and it always contains $\langle\sh,F\rangle$. This paper asks, for a natural finite family of nonlinear reversible automata, whether it contains anything else.

The family consists of the second-order elementary rules, built by the second-order technique of Fredkin \cite{Margolus1984,TM1990,Wolfram2002}; as models of statistical mechanics they were introduced by Takesue \cite{Takesue1987,Takesue1989}, who reduced them to the $88$ classes considered here \cite[Table~I]{Takesue1989}, determined their additive conserved quantities with densities on two adjacent cells \cite[Table~II]{Takesue1989}, determined, for words of length at most $6$, which classes have a local conservation law, and exhibited walls in $40$ of them \cite[Section~4, Fig.~2, Table~III]{Takesue1989} (see also \cite{HT1991,Takesue1995}). For an elementary rule $f\colon\{0,1\}^3\to\{0,1\}$ define, on $\A=\{0,1\}^2$ with cells written $(x,y)$,
\begin{equation}\label{eq:F}
F_f(c)_i=\bigl(y_i,\;x_i\oplus f(y_{i-1},y_i,y_{i+1})\bigr).
\end{equation}
Reading $y$ as the current row and $x$ as the previous one, this is the recursion $y_{t+1}=y_{t-1}\oplus f(y_t)$, reversible for every $f$, with $F_f^{-1}(c)_i=\bigl(y_i\oplus f(x_{i-1},x_i,x_{i+1}),\,x_i\bigr)$. Following Takesue \cite{Takesue1987}, one writes $f$R for $F_f$; Wolfram discusses $37$R as a reversible automaton with complex behaviour \cite[pp.~452--457]{Wolfram2002}. Reflection and the cellwise complement conjugate $F_f$ to $F_{f'}$ for the reflected and complemented rules and preserve everything we consider (Lemma~\ref{lem:classes}), so the question is one about $88$ classes, the classes of Takesue \cite[Table~I]{Takesue1989}.

\subsection{The classification}

We assign to each class one of five statuses (Definition~\ref{cl:def:status}): (G), $\Aut(X_{F_f})$ contains $\bigoplus_{m\ge1}S_m$ and hence every finite group; (Z), $\Aut(X_{F_f})=\langle\sh,F_f\rangle\cong\Z^2$; (S), $\Aut(X_{F_f})=\langle\sh,F_f\rangle\times\langle\kappa\rangle\cong\Z^2\times\Z/2$, where $\kappa$ is the cellwise complement; (A), $f$ is affine and $\Aut(X_{F_f})$ is a free abelian group strictly larger than $\langle\sh,F_f\rangle$; and (O), open.

\begin{theorem}[Classification]\label{thm:main}
The $88$ classes of second-order elementary automata, and the $256$ rules, are distributed among the five statuses as follows.
\begin{center}
\begin{tabular}{lrr}
\toprule
status & classes & rules\\
\midrule
(G) every finite group & $43$ & $104$\\
(Z) $\Aut=\langle\sh,F\rangle\cong\Z^2$ & $11$ & $40$\\
(S) $\Aut=\langle\sh,F\rangle\times\langle\kappa\rangle\cong\Z^2\times\Z/2$ & $4$ & $8$\\
(A) affine, $\Aut\cong\Z^2$ or $\Z^3$, $\supsetneq\langle\sh,F\rangle$ & $2$ & $6$\\
(O) open & $28$ & $98$\\
\midrule
total & $88$ & $256$\\
\bottomrule
\end{tabular}
\end{center}
The classes of each status are those of Tables~\ref{cl:tab:finite}, \ref{cl:tab:decided} and~\ref{cl:tab:open}. For each of the $28$ open classes, every automorphism of $X_{F_f}$ of radius at most $3$ is of the form $\sh^aF_f^b$ (computer-assisted), and $(a,b)\mapsto\sh^aF_f^b$ is an isomorphism $\Z^2\to\langle\sh,F_f\rangle$. For the class of $15$, $\Aut(X_{F_f})\cong\Z^2$ contains $\langle\sh,F_f\rangle$ with index $2$; for the class of $105$, $\Aut(X_{F_f})\cong\Z^3$; and for the class of $95$, $\Aut(X_{F_f})\cong(\mathcal N\times\mathcal N\times\Z)/\langle(\sh^{-1},\sh,-2)\rangle$, where $\mathcal N$ is the group of automorphisms of the golden mean shift that preserve the number of $1$s of every finite point \cite{PaperE}.
\end{theorem}

The theorem is proved in Section~\ref{cl:sec:main}, where the tables list the classes with the result that settles each entry. The statuses (G), (Z), (S) and (A) are quoted from the companion papers \cite{PaperA,PaperB,PaperC,PaperE} (Section~\ref{sec:imported}); what is proved here is the assembly and the statement on the open classes. The status (O) records only that none of the other four is known; its mathematical content is the census of Theorem~\ref{cl:thm:census} together with Theorem~\ref{cl:thm:A}. Whether $\Aut(X_{F_f})=\langle\sh,F_f\rangle$ for every open class is Question~\ref{q:open}. The censuses do not decide it: for the class of $95$ they contain no automorphism outside $\langle\sh,F_f\rangle$, and the group contains every finite group \eqref{imp:95}. They do show that for every open class the block maps of radius at most $4$ that commute with $F_f$ are the maps $\sh^aF^b$ and the constants, whereas for $95$ there are $20$ further ones, none of them injective, at radius $3$.

\subsection{Results used from the companion papers}\label{sec:imported}

The statuses (G), (Z), (S) and (A), and the isomorphism $\langle\sh,F_f\rangle\cong\Z^2$, rest on the following results, which are proved in four companion papers \cite{PaperA,PaperB,PaperC,PaperE}; none of them is proved here. We state them in the form in which they are used.

\begin{enumerate}
\renewcommand{\theenumi}{I\arabic{enumi}}
\renewcommand{\labelenumi}{(\theenumi)}
\item\label{imp:walls} \emph{Fixed walls} \cite{PaperC}. Let $F$ be a reversible cellular automaton of radius $1$ on $\A^{\Z}$ and $W$ a nonempty word such that $F(c)$ contains $W$ at every position at which $c$ does. If at least two symbols $s$ have the property that $W\,s\,W$ contains $W$ only at its two ends, then $\bigoplus_{m\ge1}S_m$ embeds in $\Aut(X_F)$. Of the $256$ elementary rules, $81$, in $31$ classes, have such a word of length at most $3$ (a finite check); these words are the fixed walls (i)--(viii) of Takesue \cite[Fig.~2, Table~III]{Takesue1989}.
\item\label{imp:54} \emph{Rule $54$} \cite{PaperC}. For every $m\ge1$ there is an injective homomorphism $\Psi_m\colon S_m\to\Aut(X_{F_{54}})$ whose image consists of block maps of radius at most $3m+3$, and the images for different $m$ commute; so $\bigoplus_{m\ge1}S_m$ embeds in $\Aut(X_{F_{54}})$.
\item\label{imp:133} \emph{A wall of length $4$} \cite{PaperC}. The word $0330$ is fixed by $F_{133}$ in every context (it is the wall (ix) of \cite[Table~III and p.~391]{Takesue1989}), and \eqref{imp:walls} applies to the class $\{133,161\}$, with automorphisms of radius $4$ outside $\langle\sh,F\rangle$. Every word fixed in every context by some $F_f$ lies in $\{0,3\}^k$, and among the $57$ classes with no such word of length at most $3$, the class of $133$ is the only one with such a word of any length (the second statement by a reduction to a finite search, computer-assisted in \cite{PaperC}).
\item\label{imp:hidden} \emph{Hidden-bit and time-periodic walls} \cite{PaperC}. For the classes of $59$ and $27$ (walls carrying an internal bit invisible to their neighbours) and of $127$, $179$ and $223$ (walls fixed by a power of $F$ in every context), $\bigoplus_{m\ge1}S_m$ embeds in $\Aut(X_{F_f})$. A computer search among the closures of single words under $F^{\pm1}$ (of at most $300$ words), with windows of length at most $6$ and a hidden involution satisfying the conditions (W1) and (H1) of \cite{PaperC}, in the $57$ classes without a fixed word of length at most $3$, finds such families only in these five classes \cite{PaperC}. The walls of these five classes are built from the periodic walls of \cite[Section~4]{Takesue1989}.
\item\label{imp:rot} \emph{Rotating walls} \cite{PaperC}. For $f\in\{126,219\}$, runs of cells that $F_f$ rotates through the alphabet in every context give an embedding of $\bigoplus_{m\ge1}S_m$ in $\Aut(X_{F_f})$; these runs consist of the walls (xv) and (xvi) of \cite[Fig.~2, Table~III]{Takesue1989}.
\item\label{imp:90} \emph{Markers for $90$ and $165$} \cite{PaperC}. For $f\in\{90,165\}$, $F_f$ has no blocking word (Section~\ref{sec:blocking}), and $(\bigoplus_{m\ge1}S_m)^2$ embeds in $\Aut(X_{F_f})$. The automorphisms of $X_{F_{90}}$ of radius at most $2$ are affine (computer-assisted in \cite{PaperC}; the enumeration is repeated in Table~\ref{tab:census}).
\item\label{imp:affine} \emph{Affine classes} \cite{PaperC}. For the classes of $15$, $60$, $105$ and $153$, every automorphism of $X_{F_f}$ is affine, by the isomorphism rigidity theorem of Kitchens and Schmidt \cite{KS2000}, and: for $f\in\{170,240\}$, $\Aut(X_{F_f})=\langle F,\mathrm{id}+F\rangle\cong\Z^2$, with $(\mathrm{id}+F)^2=\sh^{\pm1}F$, so that $\langle\sh,F\rangle$ has index $2$, and for $f\in\{15,85\}$ the same holds with the generator $c\mapsto c+F(c)+(1,1)^{\Z}$; for the classes of $60$ and $153$, $\Aut(X_{F_f})=\langle\sh,F\rangle\times\langle\kappa\rangle\cong\Z^2\times\Z/2$; for $f=150$, $\Aut(X_{F_f})=\langle\sh,F,\mathrm{id}+\sh F\rangle\cong\Z^3$, and for $f=105$ the same holds with $c\mapsto c+\sh F(c)+(1,1)^{\Z}$ in place of $\mathrm{id}+\sh F$. For $f=255$, $F$ adds $1$ modulo $4$ in every cell after a relabelling of the cells by $\Z/4$, $\Aut(X_{F_{255}})\cong\Z/4\times\mathcal G_4$, where $\mathcal G_4$ is the group of automorphisms $\psi$ of the full shift on $\Z/4$ with $\psi-\mathrm{id}=(\sh-\mathrm{id})J$ for a block map $J$, and $\bigoplus_{m\ge1}S_m$ embeds in $\mathcal G_4$.
\item\label{imp:perm} \emph{Permutive rules} \cite{PaperB}. If $f$ is not affine and is permutive in its left or its right variable, then $\Aut(X_{F_f})=\langle\sh,F_f\rangle\cong\Z^2$. These are the $16$ rules of the classes of $30$, $45$, $75$ and $135$.
\item\label{imp:half} \emph{Rules that read two neighbours} \cite{PaperB}. For the $8$ rules of the classes of $63$ and $187$, $\Aut(X_{F_f})=\langle\sh,F_f\rangle\cong\Z^2$.
\item\label{imp:110} \emph{Rules $110$, $14$, $46$, $79$, $24$ and $189$} \cite{PaperA}. For the $14$ rules of the classes of $110$, $14$, $46$ and $79$, $\Aut(X_{F_f})=\langle\sh,F_f\rangle\cong\Z^2$. For the $4$ rules of the classes of $24$ and $189$, $\Aut(X_{F_f})=\langle\sh,F_f\rangle\times\langle\kappa\rangle\cong\Z^2\times\Z/2$, where $\kappa$ is the cellwise complement \cite[Theorem~1.1]{PaperA}. All proofs in \cite{PaperA}, including those for $24$ and $189$, are by hand.
\item\label{imp:order} \emph{Infinite order} \cite{PaperC}. If $f$ is not a function of its middle argument alone, then $(a,b)\mapsto\sh^aF_f^b$ is injective, so $\langle\sh,F_f\rangle\cong\Z^2$ \cite[Theorem~2.4]{PaperC}. The statement is recalled as Theorem~\ref{cl:thm:A}.
\item\label{imp:94} \emph{Rule $94$} \cite{PaperE}. For $f\in\{94,122\}$, $\Aut(X_{F_f})=\langle\sh,F_f\rangle\cong\Z^2$ \cite[Theorem~A; Theorem~5.9]{PaperE}.
\item\label{imp:95} \emph{Rule $95$} \cite{PaperE}. For $f\in\{95,250\}$, $\Aut(X_{F_f})\cong(\mathcal N\times\mathcal N\times\Z)/\langle(\sh^{-1},\sh,-2)\rangle$, where $\mathcal N$ is the group of automorphisms of the golden mean shift that preserve the number of $1$s of every finite point, and $\bigoplus_{m\ge1}S_m$ embeds in $\mathcal N$, so that $\Aut(X_{F_f})$ contains every finite group; the automorphism group of the spacetime of the reversible automaton known as Rule~54, which is the system on one sublattice of a diagram of $F_{95}$, is $\mathcal N\times\mathcal N$; $F_f$ has no blocking word, and the window graphs $W^{(k,p)}_M$ of Section~\ref{sec:blocking} are strongly connected for all $(k,p)$ and $M$ \cite[Theorem~B; Theorem~6.9, Corollary~6.10, Propositions~6.17 and~6.18]{PaperE}. The smallest radius of an automorphism outside $\langle\sh,F_f\rangle$ is $7$; and the block maps of radius at most $4$ that commute with $F_{95}$, other than the maps $\sh^aF^b$ and one constant map, number $20$ at radius $3$ and $220$ at radius $4$, are induced by pairs of block maps of the golden mean shift that conserve the number of $1$s of finite points and are not both injective, and none of them is injective (computer-assisted) \cite[Proposition~6.15, Lemma~6.20 and Proposition~6.21]{PaperE}.
\item\label{imp:L} \emph{Lemma~L; rules $91$ and $155$} \cite{PaperE}. For $f\in\{91,94,95\}$, a block map of any radius that commutes with $F_f$ and fixes every fixed point of $\sh F_f$ and of $\sh^{-1}F_f$ is the identity \cite[Theorem~C; Theorem~3.1, Corollary~3.9]{PaperE}. For $91$ the group is not determined there; for $155$ the method of proof fails \cite[Proposition~8.1]{PaperE}, and the statement is verified for radius at most $6$ (computer-assisted) \cite[Remark~8.2]{PaperE}.
\end{enumerate}

The proofs of \eqref{imp:perm}--\eqref{imp:110} and of \eqref{imp:94} share a pattern: an automorphism is normalised by an element of $\langle\sh,F\rangle$ to one that fixes certain distinguished structures, and such a map is shown to be the identity. The structures are the light-speed fronts of differences for the permutive rules, and the rightmost difference between two configurations, which is stationary under $F$, for the classes of $63$ and $187$ \cite{PaperB}; the kinks among the fixed points of $\sh^{\pm1}F$ for $110$ and $46$, together with those of $F$ for $14$ and of $\sh^{\pm1}F^3$ for $79$, and for $24$ and $189$ a factor map onto rule $46$ \cite{PaperA}; and for the class of $94$ the kinks among the fixed points of $\sh^{\pm1}F$, which form a shift of finite type of positive entropy, where the scattering of kinks carries the normalisation to all these fixed points and \eqref{imp:L} concludes \cite{PaperE}. For $95$ the first step and \eqref{imp:L} hold as well, but a normalised automorphism need not be the identity; for the open class of $91$, \eqref{imp:L} gives part of the second step.

\subsection{The results of this paper}

What is proved in this paper is the following; everything else in the classification is quoted from \eqref{imp:walls}--\eqref{imp:L}.

\emph{The classification} (Theorem~\ref{thm:main}, Tables~\ref{cl:tab:finite}--\ref{cl:tab:open}). The assembly of the statuses from the imported results: which of them applies to which class, the counts, and the complete list of the permutations of the four cells that commute with $F_f$ (Remark~\ref{cl:rem:kappa}).

\emph{Censuses} (Theorems~\ref{thm:110} and~\ref{cl:thm:census}; computer-assisted). For a fixed radius $R$, commutation with $F_f$ is a finite set of constraints on the $2\cdot4^{2R+1}$ bits of the table of a block map. We enumerate all solutions with a SAT solver, re-verify each by direct evaluation, decide invertibility exactly, and repeat the enumeration with other solvers (Section~\ref{sec:method}). For the open classes and the class of $95$ the commuting block maps of radius at most $3$ are the maps $\sh^aF^b$, the constants, and, for $95$ only, $20$ non-injective maps (Theorem~\ref{cl:thm:census}). For the rules $110$ and $30$ the block maps of radius at most $4$ that commute with $F_f$ are the maps $\sh^aF^b$ and two constant maps (Theorem~\ref{thm:110}). A census at radius $4$ for the open classes, searches for walls, eigenfields and conserved quantities, and the window graphs of all classes are reported in Section~\ref{cl:sec:holdouts} (computer-assisted).

\emph{No walls in $110$R} (Proposition~\ref{cl:prop:nowalls}). Every wall mechanism of \cite{PaperC} needs a cell next to the permuted information whose trajectory does not depend on the context. Rule $110$R has no such cell, of any length or kind. The proof is by hand.

\emph{Reach of the edge-move method} (Section~\ref{cl:sec:reach}). The proof of \eqref{imp:110} works at a stuck edge, a cell beyond which a configuration agrees with a spatially periodic background and which does not move under $F^{\pm1}$. In at least $12$ classes, among them $10$ open ones and the classes of $94$ and $95$, no configuration has a stuck edge; the proof is by hand (Proposition~\ref{cl:prop:nostuck}). For $11$ further open classes a background of small period with a stuck edge has no other fixed point of $\sh^aF^b$ next to it, for small $(a,b)$ (Proposition~\ref{cl:prop:reach}, computer-assisted), and for the remaining $7$ sampled computations indicate where the method stops (Remark~\ref{cl:rem:reach}).

\subsection{Context}
The problem of finding the block maps that commute with a given one goes back to Coven, Hedlund and Rhodes \cite{CHR1979}. Voorhees \cite{Voorhees1993} describes the binary cellular automata commuting with a given one as the solutions of a system of equations in the rule table, Moore and Boykett \cite{MooreBoykett1997} show that if one of two commuting cellular automata is permutive, the other can be expressed in terms of it, and Salo and T\"orm\"a \cite{SaloTorma2013} show that the commutant of a bipermutive automaton is small and that of a bipermutive affine one consists of affine automata. Boykett \cite{Boykett2004} lists by machine, up to isomorphism, the reversible one-dimensional cellular automata with a two-cell neighbourhood whose inverse also has a two-cell neighbourhood, for up to $12$ states. The censuses of this paper enumerate, for one reversible automaton and one radius at a time, all block maps on the four-letter alphabet that commute with it.

For rule $54$, by \eqref{imp:54}, the automorphism group of the second-order automaton is not virtually $\Z^d$, while for rule $110$ it is $\Z^2$ by \eqref{imp:110}. The first-order rules $110$ and $54$ are the two elementary rules that Wolfram singles out for class~4 behaviour, $110$ as the clear example and $54$ as a possible one \cite[pp.~693--697]{Wolfram2002}; their second-order versions thus have automorphism groups of opposite kinds, and the behaviour of the first-order rule is no guide to the automorphism group of $F_f$.
In $39$ of the $43$ classes of status (G) the room comes from a conserved local structure, a wall, with an autonomous interior; walls in each of these $39$ classes were exhibited by Takesue \cite[Section~4]{Takesue1989}, and what \cite{PaperC} adds is their consequence for the automorphism group. Of the other four, three are affine, and in the class of $95$ the room comes from the automorphisms of a system of light-speed particles \cite{PaperE}. The spacetime $X_F$ has zero $\Z^2$-entropy, so Hochman's theorem that positive entropy forces every finite group into the automorphism group of a $\Z^2$ shift of finite type \cite{Hochman2010} does not apply. The techniques that prove automorphism groups small do not apply either: the complexity arguments of Cyr--Kra and Donoso--Durand--Maass--Petite \cite{CK2015,DDMP2016} are for one-dimensional subshifts whose word complexity grows at most linearly or is not superlinear, while the rows of $X_F$ are all of $\A^n$; the algebraic rigidity of Kitchens and Schmidt \cite{KS1989,KS2000}, which gives $\Z^2$ for Ledrappier's shift \cite{Ledrappier1978} (\cite{Ward1994}, \cite[Theorem~7]{BRY2018}) and determines the groups of the affine classes of $15$, $60$, $105$ and $153$ \eqref{imp:affine}, needs a group shift. Ryan-type results \cite{Ryan1972,Salo2019Ryan,Kopra2022} compute centralizers of the whole group or of specially constructed sets of two or more automata, not of one given automaton. The arguments of the companion papers are specific to the rules; whether the property $\Aut(X_F)=\langle\sh,F\rangle$ of a reversible automaton is decidable is not known (Section~\ref{sec:discussion}).

\medskip
\noindent\textbf{Organisation.} Section~\ref{sec:prelim} fixes notation and defines blocking words, window graphs and eigenfields. Section~\ref{sec:method} describes the SAT encoding, the invertibility test, the verification and the cross-check, and Section~\ref{sec:110} gives the census for $110$ and $30$. Section~\ref{sec:classification} defines the statuses, recalls Theorem~\ref{cl:thm:A}, gives the tables and the proof of Theorem~\ref{thm:main}, states the census of the open classes, proves Proposition~\ref{cl:prop:nowalls}, records what is known about the open classes, and determines the reach of the edge-move method on them. Section~\ref{sec:discussion} states open questions. The appendix describes the computations.

\section{Preliminaries}\label{sec:prelim}

Throughout, $\A=\{0,1\}^2$, and a cell $(x,y)\in\A$ is also written as the integer $2x+y\in\{0,1,2,3\}$; thus $3=(1,1)$ and $0=(0,0)$. The left shift is $\sh(c)_i=c_{i+1}$. A \emph{block map of radius $R$} is a map $g\colon\A^{\Z}\to\A^{\Z}$ of the form $g(c)_i=\gamma(c_{i-R},\dots,c_{i+R})$ for a local rule $\gamma\colon\A^{2R+1}\to\A$; we identify $g$ with the lookup table of $\gamma$, which has $4^{2R+1}$ entries of two bits each. A block map of radius $R$ is also a block map of every larger radius, so ``radius at most $R$'' and ``radius $R$'' describe the same set of maps. By Curtis--Hedlund--Lyndon, $\End(\A^{\Z})$ is the monoid of block maps and $\Aut(\A^{\Z})$ the group of invertible ones. We write
\[
\Cent_R(F)=\{g\ \text{a block map of radius}\le R:\ g\circ F=F\circ g\},\qquad \Cent(F)=\bigcup_R\Cent_R(F),
\]
for the centralizer of $F$ in $\End(\A^{\Z})$, and $\Cent^{\times}_R(F)$, $\Cent^{\times}(F)$ for the invertible elements.

For an elementary rule $f$, identified with its Wolfram number $\sum_{a,b,c}f(a,b,c)2^{4a+2b+c}$, the automaton $F_f$ is defined by \eqref{eq:F}. It has radius $1$, and $F_f^b$ has radius at most $|b|$, so $\sh^aF_f^b\in\Cent_{|a|+|b|}(F_f)$.

\begin{definition}
The \emph{spacetime subshift} of a reversible cellular automaton $F$ on $\A^{\Z}$ is
\[
X_F=\bigl\{(F^n(c))_{n\in\Z}:\ c\in\A^{\Z}\bigr\}\subset\A^{\Z^2},
\]
with the $\Z^2$-action generated by the horizontal shift $\sh$ and the vertical shift $\tau$, $(\tau z)_{n}=z_{n+1}$. An \emph{automorphism} of $X_F$ is a homeomorphism of $X_F$ commuting with both.
\end{definition}

\begin{lemma}\label{lem:spacetime}
For reversible $F$ the map $c\mapsto(F^n(c))_{n\in\Z}$ is a homeomorphism $\A^{\Z}\to X_F$ conjugating $(\sh,F)$ to $(\sh,\tau)$, and it induces an isomorphism
\[
\Aut(X_F)\;\cong\;\Cent^{\times}(F)=\{g\in\Aut(\A^{\Z}):\ gF=Fg\}.
\]
Under this isomorphism $\langle\sh,\tau\rangle$ corresponds to $\langle\sh,F\rangle$.
\end{lemma}

\begin{proof}
The map is continuous, injective (the $0$-th row recovers $c$), and surjective onto $X_F$ by definition, hence a homeomorphism of compact spaces; it carries $\sh$ to $\sh$ and $F$ to $\tau$ since $F^{n}(F(c))=F^{n+1}(c)$. A homeomorphism of $\A^{\Z}$ commuting with $\sh$ is a block map by Curtis--Hedlund--Lyndon, and commuting with $\tau$ means commuting with $F$.
\end{proof}

Thus determining $\Aut(X_F)$ means determining which invertible block maps commute with $F$, and $\Cent(F)$ is the monoid of endomorphisms of the spacetime subshift. The non-invertible elements of $\Cent(F)$ matter too: a block map that commutes with $F$ but is not injective is a non-surjective endomorphism of $X_F$, and the constant map onto a fixed point of $F$ is always one. For every rule with $f(0,0,0)=0$ and $f(1,1,1)=0$, which includes $30$, $54$, $90$ and $110$, both $(0,0)^{\Z}$ and $(1,1)^{\Z}$ are fixed by $F_f$, so $\Cent_R(F_f)$ always contains these two constants; for rule $150$, which has $f(1,1,1)=1$, only $(0,0)^{\Z}$ is fixed.


\subsection{Symmetries}
Three involutions of $\A^{\Z}$ act on second-order automata: the cellwise complement $\kappa(x,y)=(1-x,1-y)$, the reflection $\mathrm{rev}(c)_i=c_{-i}$, and the time reversal $\rho(x,y)=(y,x)$, applied in every cell. The first two are the Boolean conjugation and the reflection of \cite[Section~2]{Takesue1989}, and that exchanging the two rows reverses time is noted in \cite[p.~2500]{Takesue1987}.

\begin{lemma}\label{lem:classes}
For every elementary rule $f$, $\kappa F_f\kappa=F_{\kappa f}$ with $\kappa f(a,b,c)=f(1-a,1-b,1-c)$; $\mathrm{rev}\,F_f\,\mathrm{rev}=F_{f^r}$ with $f^r(a,b,c)=f(c,b,a)$, and $\mathrm{rev}\,\sh\,\mathrm{rev}=\sh^{-1}$; and $\rho F_f\rho=F_f^{-1}$. For $\theta\in\{\kappa,\mathrm{rev}\}$ the map $g\mapsto\theta g\theta^{-1}$ preserves the radius of block maps, maps $\Cent_R(F_f)$ onto $\Cent_R(F_{f'})$ and $\Aut(X_{F_f})$ onto $\Aut(X_{F_{f'}})$, where $f'=\kappa f$ or $f^r$, and maps $\langle\sh,F_f\rangle$ onto $\langle\sh,F_{f'}\rangle$.
\end{lemma}

\begin{proof}
By \eqref{eq:F}, $F_f(\kappa c)_i=(1-y_i,\,1-x_i\oplus f(1-y_{i-1},1-y_i,1-y_{i+1}))$, and applying $\kappa$ gives $(y_i,\,x_i\oplus\kappa f(y_{i-1},y_i,y_{i+1}))$. Similarly $\mathrm{rev}\,F_f\,\mathrm{rev}(c)_i=F_f(\mathrm{rev}\,c)_{-i}=(y_i,\,x_i\oplus f(y_{i+1},y_i,y_{i-1}))$. And $F_f(\rho c)_i=(x_i,\,y_i\oplus f(x_{i-1},x_i,x_{i+1}))$, whose image under $\rho$ is $F_f^{-1}(c)_i$. Conjugation by a cellwise map or by $\mathrm{rev}$ preserves the radius, and the remaining statements follow, since $\kappa$ commutes with $\sh$ and $\mathrm{rev}$ conjugates $\sh$ to $\sh^{-1}$.
\end{proof}

So $\Aut(X_{F_f})$, and every statement below about block maps of bounded radius, depends only on the class of $f$ under $\kappa$ and reflection. There are $88$ classes \cite{Takesue1987}, \cite[Table~I]{Takesue1989}.

\subsection{Diagrams}
For $c\in\A^{\Z}$ the \emph{diagram} of $c$ is $Y_c\colon\Z^2\to\{0,1\}$, where $Y_c(t,i)$ is the $y$-bit of $F^t(c)_i$. The $x$-row of $F(d)$ is the $y$-row of $d$, so the $x$-bit of $F^t(c)_i$ is $Y_c(t-1,i)$, and $c_i=(Y_c(-1,i),Y_c(0,i))$. By \eqref{eq:F} every diagram satisfies the \emph{spacetime relation}
\begin{equation}\label{eq:rel}
Y(t+1,i)\oplus Y(t-1,i)=f\bigl(Y(t,i-1),Y(t,i),Y(t,i+1)\bigr)\qquad\text{for all }(t,i)\in\Z^2,
\end{equation}
so two consecutive rows determine all rows, and $c\mapsto Y_c$ is a bijection from $\A^{\Z}$ onto the solutions of \eqref{eq:rel}. For $t\ge0$, $Y_c(t,i)$ depends only on $c_{[i-t,i+t]}$. Moreover $Y_{\sh c}(t,i)=Y_c(t,i+1)$, $Y_{F(c)}(t,i)=Y_c(t+1,i)$ and $Y_{\rho c}(t,i)=Y_c(-1-t,i)$, the last because $F^t(\rho c)=\rho F^{-t}(c)$.

\subsection{Blocking words, window graphs and eigenfields}\label{sec:blocking}
Let $F$ be a cellular automaton of radius $1$ on $\A^{\Z}$. A word $w$ on an interval $I$ of cells is \emph{blocking} for $F$ (cf.\ the proof of \cite[Theorem~3]{Kurka1997}; the term is from \cite{BT2000}) if there is a cell $i_0\in I$ such that $F^n(c)_{i_0}=F^n(c')_{i_0}$ for all $n\ge0$ and all configurations $c,c'$ that carry $w$ on $I$. A word that is fixed by $F$ in every context is blocking.

Let $F$ be a bijective cellular automaton on $\A^{\Z}$. For $(k,p)\in\Z^2$ with $p\ge1$ and $M\ge0$, let $W^{(k,p)}_M$ be the graph on $\A^{M+1}$ with an edge from each word to each of its possible images under $\sh^kF^p$, the cells outside the window ranging over all values: there is an edge from $c_{[0,M]}$ to $\sh^kF^p(c)_{[0,M]}$ for every $c\in\A^{\Z}$. The graph $W^{(1,1)}_M$ is the light-speed window graph $\Lambda_M$ used for the permutive rules \cite{PaperB}, and $W^{(0,1)}_M$ is a zero-speed graph. We write $(\Lambda^{\pm})$ for the condition that $W^{(1,1)}_M$ and $W^{(-1,1)}_M$ are strongly connected for every $M$, and we call the condition that $W^{(0,1)}_M$ is strongly connected for every $M$ the \emph{zero-speed condition}. A nonconstant block map $\vartheta\colon\A^{\Z}\to B^{\Z}$ onto sequences over a finite alphabet $B$ with $\vartheta\circ F^p=\sh^{-k}\circ\vartheta$ is an \emph{eigenfield of speed $k/p$}, and a \emph{frozen field} if $k=0$; composing with the indicator of a letter shows that one exists if and only if one exists with $B=\{0,1\}$.

\begin{lemma}[{\cite[Lemma~2.5]{PaperE}}]\label{cl:lem:eigen}
Let $F$ be a bijective cellular automaton on $\A^{\Z}$ and $(k,p)\in\Z^2$ with $p\ge1$. The following are equivalent.
\begin{enumerate}
\item[(1)] The graph $W^{(k,p)}_M$ is strongly connected for every $M\ge0$.
\item[(2)] There is no nonconstant block map $\vartheta\colon\A^{\Z}\to\{0,1\}^{\Z}$ with $\vartheta\circ F^p=\sh^{-k}\circ\vartheta$.
\item[(3)] Every continuous function $v\colon\A^{\Z}\to\{0,1\}$ with $v\circ\sh^kF^p=v$ is constant.
\end{enumerate}
\end{lemma}

By the lemma, $(\Lambda^{\pm})$ says that $F$ has no eigenfield of speed $\pm1$, and the zero-speed condition says that $F$ has no frozen field. A frozen field is a local conservation law in the sense of Takesue \cite{Takesue1987}, \cite[Section~4]{Takesue1989}, and $W^{(0,1)}_M$ is the graph of the transition matrix on words of $M+1$ cells considered there; the equivalence of (1) and (2) for $(k,p)=(0,1)$ is stated in \cite[Section~4]{Takesue1989} for the second-order elementary automata (see also \cite[Proposition~4.1]{HT1991}).

\section{The computation}\label{sec:method}

\subsection{Commutation as a finite constraint}
Let $g$ be a block map of radius $R$ with local rule $\gamma$, and $F=F_f$. Both $g\circ F$ and $F\circ g$ are block maps of radius $R+1$, so they are equal if and only if their local rules agree on every window $w\in\A^{2R+3}$. Write $n=2R+1$ and, for a window $w=(w_0,\dots,w_{2R+2})$, let $e_1,e_2,e_3$ be the three subwords of length $n$ starting at positions $0,1,2$ and let $e_L=F(w)$ be the word of length $n$ obtained by applying the local rule of $F$ to $w$. Then
\[
(g\circ F)(w)=\gamma(e_L),\qquad (F\circ g)(w)=\Bigl(\gamma(e_2)_y,\ \gamma(e_2)_x\oplus f\bigl(\gamma(e_1)_y,\gamma(e_2)_y,\gamma(e_3)_y\bigr)\Bigr),
\]
where $\gamma(e)_x,\gamma(e)_y$ are the two output bits. Writing $X_e,Y_e$ for the two unknown bits of table entry $e$, each window contributes the two constraints
\begin{equation}\label{eq:constraints}
X_{e_L}=Y_{e_2},\qquad Y_{e_L}\oplus X_{e_2}=f(Y_{e_1},Y_{e_2},Y_{e_3}).
\end{equation}
The first is two binary clauses; the second is a constraint on at most five variables, encoded by its truth table in at most sixteen clauses (the indices $e_L,e_1,e_2,e_3$ may coincide, and the encoding is over the set of distinct variables). Altogether there are $2\cdot4^{2R+1}$ variables and at most $18\cdot4^{2R+3}$ clauses; for rule $110$ there are $17{,}680$ clauses at $R=1$, $4.7\times10^6$ at $R=3$ and $7.5\times10^7$ at $R=4$.

\subsection{Enumeration and verification}
All solutions are enumerated by repeated solving with blocking clauses, using PySAT \cite{PySAT} with CaDiCaL~1.5.3 \cite{CaDiCaL} for $R\le3$ and MiniSat~2.2 \cite{MiniSat} at $R=4$, where CaDiCaL's memory use exceeded the machine. Every table returned by the solver is re-checked, independently of the CNF, by evaluating $g\circ F$ and $F\circ g$ on all $4^{2R+3}$ windows in NumPy. Each solution is then classified: it is matched against the $\sh^aF^b$ with $|a|+|b|\le R$, which are built directly by composing local rules; the two constants are recognised; everything else is labelled \emph{other}. The program asserts that every $\sh^aF^b$ with $|a|+|b|\le R$ occurs among the solutions, which checks the encoding against the only maps known in advance to commute.

Every solution is thus verified independently of the solver. Exhaustiveness is of another kind: the clauses are necessary conditions for commutation, so every commuting block map is a solution, and that there is no solution beyond those found rests on the final answer of the solver that the clauses, together with the blocking clauses of the solutions found, are unsatisfiable. This answer was cross-checked in two ways (\texttt{code/census/crosscheck.py}, Appendix~\ref{app:census}). First, every enumeration was repeated with other solvers: for $R\le3$ with Glucose~4 \cite{Glucose}, MiniSat~2.2 and once more CaDiCaL~1.5.3, on all $200$ instances that finish below the cap; for $R=4$ with Glucose~4 and CaDiCaL~1.5.3 for the rules $110$ and $30$, and with Glucose~4 for the other $39$ rules that finish below the cap. In every instance the solvers return the same set of tables, and this set agrees with the log of the enumeration. Second, for $R\le3$ the unsatisfiable formula was given to Glucose~4 with proof logging, and the clausal proof it emits was checked by a short forward reverse-unit-propagation checker that uses no SAT solver; this was done for $197$ of the $200$ instances, among them all those behind Theorem~\ref{cl:thm:census} and the cases $R\le3$ of Theorem~\ref{thm:110}, the exceptions being three instances with more than $1700$ solutions. No proof was checked at $R=4$, where exhaustiveness rests on the agreement of two or three solvers.

\subsection{Invertibility}
A block map is an automorphism of $\A^{\Z}$ if and only if it is injective. Injectivity of a block map of radius $R$ is decided exactly by the pair graph: its vertices are the pairs $(u,v)$ of words of length $2R$, with an edge $(u,v)\to(u',v')$ labelled $(a,b)\in\A^2$ whenever $u'=ua$ and $v'=vb$ after dropping the first letter, and $\gamma(ua)=\gamma(vb)$. Then $g$ fails to be injective if and only if a bi-infinite path passes through an off-diagonal vertex, which is decided by iteratively deleting vertices without a predecessor or without a successor. We implemented this twice, in pure Python and vectorised, and asserted agreement at $R=1$; at $R\le2$ every map found injective was also checked to be injective on all spatially periodic configurations of period at most $2R+4$.

For $R=2$ with rule $54$, where the full centralizer is too large to enumerate, we enumerated only the invertible part, as follows. An invertible block map is surjective, and a surjective block map of radius $R$ is \emph{balanced} \cite[Theorem~5.4]{Hedlund1969}: each output symbol occurs exactly $4^{2R}$ times in the table. Balance is added to the CNF as four cardinality constraints, the balanced commuting maps are enumerated, and each is submitted to the pair-graph test. Balance is necessary and the pair-graph test is exact, so this yields $\Cent^{\times}_R(F)$ exactly.

\subsection{Linear rules}
When $f$ is linear over $\F_2$, the constraints \eqref{eq:constraints} are homogeneous linear equations in the table bits and $\Cent_R(F)$ is an $\F_2$-subspace of table space. Its dimension was computed by Gaussian elimination: the equalities $X_{e_L}=Y_{e_2}$ are merged by union--find, which reduces the $2\cdot4^{2R+1}$ bits to $4^{2R}$ classes, and the remaining equations are reduced by bit-packed elimination. For rule $90$ and $R\le2$ the result was compared with the SAT enumeration of all $2^{\dim}$ solutions; this tests the encoding and the solver on a control whose solution set is known independently.


\section{Rules 110 and 30}\label{sec:110}

\begin{table}[t]
\caption{Block maps of radius at most $R$ commuting with $F_f$. ``Named'' are the maps $\sh^aF^b$ with $|a|+|b|\le R$, of which there are $2R^2+2R+1$. ``Other'' are maps that are neither named nor constant. The column ``invertible'' counts the invertible commuting maps, that is, the automorphisms of $X_F$ of radius at most $R$. The rows for $54$R, $90$R and $150$R are controls; they repeat enumerations of \cite{PaperC} (computer-assisted there), and the numbers of invertible maps for $150$R are those of \cite{PaperC}. $D$ and the affine maps are described after the proof of Theorem~\ref{thm:110}, and ``linear commutant'' means that the commuting maps are the linear maps $(aI+bM)c$.}
\label{tab:census}
\centering\small
\begin{tabular}{llrrrrrl}
\toprule
rule & $R$ & commuting & named & constants & other & invertible & automorphisms outside $\langle\sh,F\rangle$\\
\midrule
110R & 1 & 7 & 5 & 2 & 0 & 5 & none\\
110R & 2 & 15 & 13 & 2 & 0 & 13 & none\\
110R & 3 & 27 & 25 & 2 & 0 & 25 & none\\
110R & 4 & 43 & 41 & 2 & 0 & 41 & none\\
30R & 1 & 7 & 5 & 2 & 0 & 5 & none\\
30R & 2 & 15 & 13 & 2 & 0 & 13 & none\\
30R & 3 & 27 & 25 & 2 & 0 & 25 & none\\
30R & 4 & 43 & 41 & 2 & 0 & 41 & none\\
54R & 1 & 13 & 5 & 2 & 6 & 5 & none (the $6$ others are not injective)\\
54R & 2 & $>2\times10^4$ & 13 & 2 & rest & 23 & $10$: $F^bD^k$, $|b|\le2$, $k=1,2$\\
90R & 1 & $2^{5}$ & 5 & 2 & 25 & 10 & $5$, all affine\\
90R & 2 & $2^{17}$ & 13 & 2 & rest & 26 & $13$, all affine\\
90R & 3 & $2^{65}$ & 25 & 2 & rest & --- & ---\\
150R & 1 & $2^{4}$ & 5 & 1 & 10 & 9 & $4$; linear commutant\\
150R & 2 & $2^{8}$ & 13 & 1 & rest & 43 & $30$; linear commutant\\
150R & 3 & $2^{12}$ & 25 & 1 & rest & 115 & $90$; linear commutant\\
\bottomrule
\end{tabular}
\end{table}

\begin{theorem}[computer-assisted]\label{thm:110}
Let $f\in\{110,30\}$ and $F=F_f$. For every $R\le 4$, the block maps of radius at most $R$ on $(\{0,1\}^2)^{\Z}$ that commute with $F$ are exactly the maps $\sh^aF^b$ with $|a|+|b|\le R$ and the two constant maps onto the fixed configurations $(0,0)^{\Z}$ and $(1,1)^{\Z}$.
\end{theorem}

\begin{proof}
The enumeration described in Section~\ref{sec:method} was run for $f=110$ and $f=30$ at $R=1,2,3,4$. In every case the solver returned exactly $2R^2+2R+3$ tables, each re-verified to commute with $F$ by direct evaluation; $2R^2+2R+1$ of them coincided with the tables of $\sh^aF^b$, $|a|+|b|\le R$, and the remaining two were the constant maps onto $(0,0)^{\Z}$ and $(1,1)^{\Z}$. The $\sh^aF^b$ are invertible and the constants are not. The counts are in Table~\ref{tab:census}; the radius-$4$ instances have $524{,}288$ variables and $7.5\times10^7$ clauses. The enumerations were repeated with other solvers, with identical solution sets (Section~\ref{sec:method}).
\end{proof}

The theorem says that $\Cent_4(F)$, the set of endomorphisms of $X_F$ of radius at most $4$, consists of the elements of $\langle\sh,F\rangle$ of radius at most $4$ and the two constants. For automorphisms the conclusion holds at every radius: $\Aut(X_{F_{110}})=\langle\sh,F_{110}\rangle$ by \eqref{imp:110}, and $\Aut(X_{F_{30}})=\langle\sh,F_{30}\rangle$ by \eqref{imp:perm}; for the non-invertible commuting block maps of these two rules, Theorem~\ref{thm:110} is what is known. It is a statement about all block maps on the four-symbol alphabet: a search among the maps that act on the two layers separately, $g(x,y)=(\phi x,\phi y)$ for a binary block map $\phi$, would see only the shifts and the two constants, since $\sh^aF^b$ mixes the layers for $b\ne0$.

The rows of Table~\ref{tab:census} for $54$, $90$ and $150$ serve as controls of the encoding; the automorphisms they find are explained in \cite{PaperC}: for $54$R, $D$ is the relabelling of radius $2$ that acts as the $3$-cycle $0\to1\to2\to0$ on every single cell between two walls $33$ and as the identity elsewhere, and for $90$R the automorphisms of radius at most $2$ outside $\langle\sh,F\rangle$ are affine maps $c\mapsto(aI+bM)c+k$ in the notation of Remark~\ref{cl:rem:affine}.

\section{The classification of second-order elementary automata}\label{sec:classification}

By Lemma~\ref{lem:classes}, $\Aut(X_{F_f})$ and every statement below about block maps of bounded radius depend only on the class of $f$ under reflection and the cellwise complement. This section states what each status means, lists the classes with the results that settle them, proves Theorem~\ref{thm:main}, and collects what the computations say about the open classes.

\subsection{What the entries mean}\label{cl:sec:status}

\begin{definition}\label{cl:def:status}
A class has one of the following five statuses.
\begin{enumerate}
\item[(G)] \emph{Every finite group.} There is an injective homomorphism $\bigoplus_{m\ge1}S_m\to\Aut(X_{F_f})$. Then $\Aut(X_{F_f})$ contains every finite group and is not virtually $\Z^n$ for any $n$. This is a lower bound. Only for the classes of $255$ and $95$ is the whole group described: for $255$ as $\Z/4\times\mathcal G_4$ \eqref{imp:affine}, and for $95$ in terms of the group $\mathcal N$ \eqref{imp:95}.
\item[(Z)] \emph{$\Z^2$.} $\Aut(X_{F_f})=\langle\sh,F_f\rangle$, and $(a,b)\mapsto\sh^aF_f^b$ is an isomorphism $\Z^2\to\langle\sh,F_f\rangle$.
\item[(S)] \emph{$\Z^2\times\Z/2$.} $\Aut(X_{F_f})=\langle\sh,F_f\rangle\times\langle\kappa\rangle\cong\Z^2\times\Z/2$, where $\kappa$ is the cellwise complement.
\item[(A)] \emph{Affine, free abelian.} $f$ is affine, and $\Aut(X_{F_f})$ is a free abelian group of finite rank that contains $\langle\sh,F_f\rangle$ as a proper subgroup.
\item[(O)] \emph{Open.} None of the above is known.
\end{enumerate}
\end{definition}

The statuses (G), (Z), (S) and (A) exclude one another: a group of status (Z), (S) or (A) is virtually $\Z^n$, a group of status (A) is torsion-free and larger than $\langle\sh,F_f\rangle$. The status (O) is a statement about what is known, not about the automata; its mathematical content is the census of Theorem~\ref{cl:thm:census}, which finds no automorphism outside $\langle\sh,F_f\rangle$ of radius at most $3$. The census says nothing about larger radii.

\begin{remark}\label{cl:rem:affine}
For an affine rule $f(a,b,c)=\alpha a\oplus\beta b\oplus\gamma c\oplus\varepsilon$, let $\mathcal R=\F_2[s,s^{-1}]$ act on $\F_2^{\Z}$ through $s=\sh$. Then $F_f(c)=Mc+k_0$, where $M=\begin{psmallmatrix}0&1\\1&p\end{psmallmatrix}$ with $p=\alpha s^{-1}+\beta+\gamma s$ acts on columns $(x,y)^{\mathsf T}$ and $k_0=(0,\varepsilon)^{\Z}$. An affine block map $c\mapsto Nc+k$, with $N$ linear and $k$ a constant configuration, commutes with $F_f$ if and only if $NM=MN$ and $(N+I)k_0=(M+I)k$, and the matrices commuting with $M$ are the $aI+bM$ with $a,b\in\mathcal R$. There are nine affine classes. For those of $15$, $60$, $105$ and $153$ every automorphism is of this form and the groups are determined \eqref{imp:affine}; those of $0$, $51$, $90$, $165$ and $255$ have status (G). For $255$, $f\equiv1$ and $F_{255}(x,y)=(y,\neg x)$ is the cellwise permutation $0\mapsto1\mapsto3\mapsto2\mapsto0$, so $F_{255}^2=\kappa$, $F_{255}^4=\mathrm{id}$ and $\langle\sh,F_{255}\rangle\cong\Z\times\Z/4$.
\end{remark}

\subsection{The group generated by the shift and the automaton}\label{cl:sec:Z2}

For an elementary rule $f$ let $\Phi_f\subseteq\{0,1\}^{\Z}$ be the set of sequences $\varphi$ with
\[
\varphi(k-1)\oplus\varphi(k+1)=f\bigl(\varphi(k-1),\varphi(k),\varphi(k+1)\bigr)\qquad\text{for all }k\in\Z.
\]

\begin{lemma}[{\cite[Lemma~2.3]{PaperA}}]\label{cl:lem:Z2}
The maps $\varphi\mapsto c$, $c_i=(\varphi(i+1),\varphi(i))$, and $\varphi\mapsto d$, $d_i=(\varphi(i-1),\varphi(i))$, are bijections from $\Phi_f$ onto the fixed points of $\sh F_f$ and of $\sh^{-1}F_f$ respectively.
\end{lemma}

The set $\Phi_f$ is the shift of finite type on $\{0,1\}$ given by the allowed triples $(u,v,w)$ with $u\oplus w=f(u,v,w)$. It has a point that is not periodic if and only if its graph on pairs has a strongly connected component that is not a single cycle, or a path between two distinct cycles; this is a finite check, and its outcome is recorded in Tables~\ref{cl:tab:decided} and~\ref{cl:tab:open}. For $110$ the allowed triples are $000,001,011,110,111$, and $\Phi_{110}$ consists of $0^{\Z}$, $1^{\Z}$ and the steps $0^{(-\infty,p]}1^{(p,\infty)}$. We call a point of $\mathrm{Fix}(\sh^aF_f^b)$ that is not spatially periodic but agrees with spatially periodic configurations on a left and on a right half-line a \emph{kink}; the fixed points of $\sh^{\pm1}F_{110}$ that correspond to the steps are the kinks on which the proof of \eqref{imp:110} is built \cite{PaperA}. If $\Phi_f$ contains a sequence that is not periodic, then the corresponding $c$ and $d$ are not periodic under $\sh$, and $\sh^aF_f^b=\mathrm{id}$ gives $c=\sh^{a-b}c$ and $d=\sh^{a+b}d$, so $a=b=0$ and $\langle\sh,F_f\rangle\cong\Z^2$. The following theorem gives this without a hypothesis on $\Phi_f$.

A block map $g$ \emph{depends on} a bit of the cell $c_i$ (for the output bit $y(g(c)_0)$, say) if two configurations that differ only in that bit have different outputs. For a nonconstant block map $g$ let $r^+(g)$ be the largest and $-r^-(g)$ the smallest $i$ such that $g(c)_0$ depends on $c_i$.

\begin{theorem}[{\cite[Theorem~2.4]{PaperC}}]\label{cl:thm:A}
Let $f$ be an elementary rule, $F=F_f$ and $n\ge1$.
\begin{enumerate}
\item If $f$ depends on its third argument, then $y(F^n(c)_0)$ depends on $y(c_n)$, and $x(F^{-n}(c)_0)$ depends on $x(c_n)$. So $r^+(F^n)=r^+(F^{-n})=n$.
\item If $f$ depends on its first argument, then $y(F^n(c)_0)$ depends on $y(c_{-n})$, and $x(F^{-n}(c)_0)$ depends on $x(c_{-n})$. So $r^-(F^n)=r^-(F^{-n})=n$.
\item The map $(a,b)\mapsto\sh^aF^b$ is an isomorphism $\Z^2\to\langle\sh,F\rangle$ if and only if $f$ is not a function of its middle argument alone.
\end{enumerate}
\end{theorem}

The rules that are functions of their middle argument alone are $0$, $51$, $204$ and $255$; for them $F$ acts by the same permutation of $\A$ in every cell, of order $2$, $3$, $3$ and $4$ respectively. So $F_f$ has infinite order and $\langle\sh,F_f\rangle\cong\Z^2$ in every class except the three classes $\{0\}$, $\{51,204\}$ and $\{255\}$. Among the classes of statuses (Z), (S), (A) and (O), those in which $\Phi_f$ has only periodic points, so that the theorem is needed, are the $21$ open classes with finite $\Phi_f$ (Section~\ref{cl:sec:holdouts}) and the classes of $45$, $135$, $63$, $187$, $189$, $60$, $153$, $15$ and $105$.

\begin{remark}[cellwise symmetries]\label{cl:rem:kappa}
Let $f$ be none of $0$, $51$, $204$, $255$. By Theorem~\ref{cl:thm:A} the group $\langle\sh,F_f\rangle\cong\Z^2$ is torsion-free, so a permutation $\pi\ne\mathrm{id}$ of the four cells that commutes with $F_f$ when applied in every cell does not lie in it, and $\langle\sh,F_f,\pi\rangle=\langle\sh,F_f\rangle\times\langle\pi\rangle$. A check over the $23$ nontrivial permutations of $\A$ (\texttt{code/census/classes88.py}) gives the complete list of such symmetries. The complement $\kappa$ is the only one for the ten classes of $24$, $36$, $60$, $90$, $126$, $129$, $153$, $165$, $189$ and $219$, and none of the other classes apart from those of $0$, $51$ and $255$ has any. These ten classes and those of $0$ and $255$ are the self-conjugate classes of \cite[(A.4)]{Takesue1989}; what the check adds is that no other cell permutation commutes. The six pairs listed there consist of a self-conjugate class and the class of its complement, and the correspondence between their orbits \cite[(A.1)--(A.3)]{Takesue1989} implies $F_{1\oplus f}^4=F_f^4$ for a self-conjugate rule $f$; for the pair of $24$ and $189$ this is the relation $F_{189}^4=F_{66}^4$ of Table~\ref{cl:tab:decided}. For $126$ and $219$ the resulting inclusion $\Aut(X_{F_f})\supseteq\Z^2\times\Z/2$ is \cite[Proposition~2.5]{PaperC}; for $24$ and $189$, and for $60$ and $153$, equality holds, by \cite[Theorem~1.1(b)]{PaperA} and \eqref{imp:affine}, and these are the four classes of status (S). The other six of the ten classes have status (G): those of $36$ and $129$ have fixed walls, those of $126$ and $219$ rotating walls, with automorphisms of radius $2$ outside $\langle\sh,F_f\rangle\times\langle\kappa\rangle$ \cite{PaperC}, and those of $90$ and $165$ are settled by \eqref{imp:90}. For $0$, $F_0$ is the transposition $(1\,2)$ of the cells, and the commuting cell permutations are $F_0$, $(0\,3)$ and $\kappa=(0\,3)(1\,2)$; for $51$, $F_{51}$ is the $3$-cycle $(0\,1\,2)$, and the commuting cell permutations are $F_{51}^{\pm1}$; for $255$ they are $F_{255}$, $F_{255}^2=\kappa$ and $F_{255}^3$.
\end{remark}

\subsection{The tables and the proof of Theorem~\ref{thm:main}}\label{cl:sec:main}

The classes are listed with their statuses in Tables~\ref{cl:tab:finite}, \ref{cl:tab:decided} and~\ref{cl:tab:open}. The counts of Theorem~\ref{thm:main} are composed as follows. In status (G), $31$ classes ($81$ rules) have a word of length at most $3$ fixed in every context; the class of $133$ ($2$ rules) has one of length $4$; $5$ classes ($4+4+2+2+2$ rules), those of $59$, $27$, $127$, $179$ and $223$, have hidden-bit or time-periodic walls; the classes of $126$ and $219$ ($1+1$ rules) have rotating walls (in these $39$ classes the walls are those of Takesue \cite[Section~4]{Takesue1989} or are built from them, and the embeddings are those of \cite{PaperC}); the affine classes of $90$ and $165$ ($1+1$ rules) have no wall of any of these kinds and get $(\bigoplus_mS_m)^2$ from marker permutations; the affine class of $255$ ($1$ rule) gets $\bigoplus_mS_m$ from marker permutations of a difference sequence; and the class of $95$ ($2$ rules), which is not affine, gets $\bigoplus_mS_m$ from marker permutations of its light-speed particles. This gives $43$ classes and $104$ rules. In status (Z), the $4$ outer-permutive classes have $16$ rules, the classes of $63$ and $187$ have $4+4$, those of $110$, $14$, $46$ and $79$ have $4+4+2+4$, and the class of $94$ has $2$: $11$ classes and $40$ rules. The $8$ rules of status (S) are $24$, $66$, $189$, $231$, $60$, $102$, $153$ and $195$, and the $6$ rules of status (A) are $15$, $85$, $170$, $240$, $105$ and $150$.

\begin{proof}[Proof of Theorem~\ref{thm:main}]
The classes under reflection and the cellwise complement, the affine and the outer-permutive rules, and the words of length at most $3$ fixed in every context are computed by \texttt{code/census/classes88.py}. (G) For the $31$ classes with a fixed word of length at most $3$, and for the class of $133$, \eqref{imp:walls} and \eqref{imp:133}; for $59$, $27$, $127$, $179$ and $223$, \eqref{imp:hidden}; for $219$ and $126$, \eqref{imp:rot}; for $90$ and $165$, \eqref{imp:90}; for $255$, \eqref{imp:affine}; for $95$, \eqref{imp:95}. (Z) \eqref{imp:perm} for the classes of $30$, $45$, $75$ and $135$; \eqref{imp:half} for the classes of $63$ and $187$; \eqref{imp:110} for the classes of $110$, $14$, $46$ and $79$; \eqref{imp:94} for the class of $94$. (S) \eqref{imp:110} for $24$ and $189$; \eqref{imp:affine} for $60$ and $153$. (A) \eqref{imp:affine} for $15$ and $105$. The remaining $28$ classes are those of Table~\ref{cl:tab:open}; for them the statement on automorphisms of radius at most $3$ is Theorem~\ref{cl:thm:census}, and $\langle\sh,F_f\rangle\cong\Z^2$ is Theorem~\ref{cl:thm:A}, since none of them contains $0$, $51$, $204$ or $255$. The descriptions of the groups in the last sentence of the theorem are those of \eqref{imp:110}, \eqref{imp:affine} and \eqref{imp:95}.

Three remarks on the assembly. The affine classes of $0$ and $51$ have fixed walls and are listed under (G). The affine classes of $90$ and $165$ also have $\kappa$ as their only nontrivial cellwise symmetry (Remark~\ref{cl:rem:kappa}); they are listed under (G) by \eqref{imp:90}. No class of status (Z) is affine; none can have a wall of the kinds of \cite{PaperC}, since such a wall gives status (G).
\end{proof}

\begin{table}[!htbp]
\caption{Status (G): $\bigoplus_mS_m$ embeds in $\Aut(X_{F_f})$. The fixed walls are the walls (i)--(ix) of Takesue \cite[Table~III]{Takesue1989}, and the other walls are built from his periodic walls. For a fixed wall, $W$ is a shortest word fixed by $F_f$ in every context (cells written $0$--$3$); for hidden-bit walls, $k$ is the length of the wall words; for time-periodic walls, each wall word is fixed by $F^P$ in every context; rotating walls are runs of cells that $F$ maps by the $4$-cycle $\pi$ in every context; for $90$ and $165$ the certificate is the marker construction of \eqref{imp:90}, and for $255$ that of \eqref{imp:affine}; these three classes are also affine; for $95$ the certificate is the marker construction of \cite[Lemma~6.8]{PaperE} on the right-moving particles of the two sublattices. The last column gives the result that settles the entry; all but the entry for $95$, which is \cite[Corollary~6.10]{PaperE}, are proved in \cite{PaperC}. The representative is the smallest member.}
\label{cl:tab:finite}
\centering\small\renewcommand{\arraystretch}{0.9}
\begin{tabular}{rlll}
\toprule
rep. & members & certificate & theorem\\
\midrule
0 & 0 & fixed wall $W=0$ & \cite{PaperC}, \eqref{imp:walls}\\
1 & 1, 128 & fixed wall $W=3$ & \cite{PaperC}, \eqref{imp:walls}\\
2 & 2, 8, 16, 64 & fixed wall $W=3$ & \cite{PaperC}, \eqref{imp:walls}\\
3 & 3, 17, 136, 192 & fixed wall $W=3$ & \cite{PaperC}, \eqref{imp:walls}\\
4 & 4, 32 & fixed wall $W=0$ & \cite{PaperC}, \eqref{imp:walls}\\
5 & 5, 160 & fixed wall $W=33$ & \cite{PaperC}, \eqref{imp:walls}\\
6 & 6, 20, 40, 96 & fixed wall $W=33$ & \cite{PaperC}, \eqref{imp:walls}\\
7 & 7, 21, 168, 224 & fixed wall $W=33$ & \cite{PaperC}, \eqref{imp:walls}\\
9 & 9, 65, 130, 144 & fixed wall $W=30$ & \cite{PaperC}, \eqref{imp:walls}\\
10 & 10, 80 & fixed wall $W=30$ & \cite{PaperC}, \eqref{imp:walls}\\
11 & 11, 81, 138, 208 & fixed wall $W=30$ & \cite{PaperC}, \eqref{imp:walls}\\
12 & 12, 34, 48, 68 & fixed wall $W=0$ & \cite{PaperC}, \eqref{imp:walls}\\
18 & 18, 72 & fixed wall $W=3$ & \cite{PaperC}, \eqref{imp:walls}\\
19 & 19, 200 & fixed wall $W=3$ & \cite{PaperC}, \eqref{imp:walls}\\
22 & 22, 104 & fixed wall $W=33$ & \cite{PaperC}, \eqref{imp:walls}\\
23 & 23, 232 & fixed wall $W=33$ & \cite{PaperC}, \eqref{imp:walls}\\
27 & 27, 83, 202, 216 & hidden-bit walls, $k=6$ & \cite{PaperC}, \eqref{imp:hidden}\\
33 & 33, 132 & fixed wall $W=3$ & \cite{PaperC}, \eqref{imp:walls}\\
35 & 35, 49, 140, 196 & fixed wall $W=3$ & \cite{PaperC}, \eqref{imp:walls}\\
36 & 36 & fixed wall $W=00$ & \cite{PaperC}, \eqref{imp:walls}\\
37 & 37, 164 & fixed wall $W=33$ & \cite{PaperC}, \eqref{imp:walls}\\
38 & 38, 44, 52, 100 & fixed wall $W=33$ & \cite{PaperC}, \eqref{imp:walls}\\
39 & 39, 53, 172, 228 & fixed wall $W=33$ & \cite{PaperC}, \eqref{imp:walls}\\
50 & 50, 76 & fixed wall $W=3$ & \cite{PaperC}, \eqref{imp:walls}\\
51 & 51, 204 & fixed wall $W=3$ & \cite{PaperC}, \eqref{imp:walls}\\
54 & 54, 108 & fixed wall $W=33$ & \cite{PaperC}, \eqref{imp:walls}, \eqref{imp:54}\\
55 & 55, 236 & fixed wall $W=33$ & \cite{PaperC}, \eqref{imp:walls}\\
59 & 59, 115, 206, 220 & hidden-bit walls, $k=4$ & \cite{PaperC}, \eqref{imp:hidden}\\
73 & 73, 146 & fixed wall $W=030$ & \cite{PaperC}, \eqref{imp:walls}\\
90 & 90 & affine; markers, $(\bigoplus_mS_m)^2$ & \cite{PaperC}, \eqref{imp:90}\\
95 & 95, 250 & light-speed particles; markers on one family & \cite{PaperE}, \eqref{imp:95}\\
126 & 126 & rotating walls, $c_{i+1}=\kappa(c_i)$ & \cite{PaperC}, \eqref{imp:rot}\\
127 & 127, 254 & time-periodic walls, $P=8$ & \cite{PaperC}, \eqref{imp:hidden}\\
129 & 129 & fixed wall $W=03$ & \cite{PaperC}, \eqref{imp:walls}\\
131 & 131, 137, 145, 193 & fixed wall $W=30$ & \cite{PaperC}, \eqref{imp:walls}\\
133 & 133, 161 & fixed wall $W=0330$ ($k=4$) & \cite{PaperC}, \eqref{imp:133}\\
139 & 139, 209 & fixed wall $W=30$ & \cite{PaperC}, \eqref{imp:walls}\\
147 & 147, 201 & fixed wall $W=303$ & \cite{PaperC}, \eqref{imp:walls}\\
165 & 165 & affine; markers, $(\bigoplus_mS_m)^2$ & \cite{PaperC}, \eqref{imp:90}\\
179 & 179, 205 & time-periodic walls, $P=6$ & \cite{PaperC}, \eqref{imp:hidden}\\
219 & 219 & rotating walls, $c_{i+1}=c_i$ & \cite{PaperC}, \eqref{imp:rot}\\
223 & 223, 251 & time-periodic walls, $P=8$ & \cite{PaperC}, \eqref{imp:hidden}\\
255 & 255 & affine; markers on $x_{i+1}-x_i\in\Z/4$ & \cite{PaperC}, \eqref{imp:affine}\\
\bottomrule
\end{tabular}
\end{table}

\begin{table}[!htbp]
\caption{Statuses (Z), (S) and (A). In the certificate column, ``$\Phi_f$'' means that $\Phi_f$ has a non-periodic point (Section~\ref{cl:sec:Z2}); $\langle\sh,F_f\rangle\cong\Z^2$ holds for every class in the table by Theorem~\ref{cl:thm:A}. The last column gives the result that settles the entry.}
\label{cl:tab:decided}
\centering\footnotesize\renewcommand{\arraystretch}{0.9}
\begin{tabular}{rllll}
\toprule
rep. & members & status & certificate & theorem\\
\midrule
30 & 30, 86, 106, 120 & (Z) & outer-permutive, not affine; $\Phi_f$ & \cite{PaperB}, \eqref{imp:perm}\\
45 & 45, 101, 166, 180 & (Z) & outer-permutive, not affine & \cite{PaperB}, \eqref{imp:perm}\\
75 & 75, 89, 154, 210 & (Z) & outer-permutive, not affine; $\Phi_f$ & \cite{PaperB}, \eqref{imp:perm}\\
135 & 135, 149, 169, 225 & (Z) & outer-permutive, not affine & \cite{PaperB}, \eqref{imp:perm}\\
63 & 63, 119, 238, 252 & (Z) & two neighbours; stationary difference & \cite{PaperB}, \eqref{imp:half}\\
187 & 187, 207, 221, 243 & (Z) & two neighbours; stationary difference & \cite{PaperB}, \eqref{imp:half}\\
62 (class of $110$) & 62, 110, 118, 124 & (Z) & kinks; edge moves; $\Phi_f$ & \cite{PaperA}, \eqref{imp:110}\\
14 & 14, 42, 84, 112 & (Z) & kinks, static kinks; edge moves; $\Phi_f$ & \cite{PaperA}, \eqref{imp:110}\\
46 & 46, 116 & (Z) & kinks; bond bits; $\Phi_f$ & \cite{PaperA}, \eqref{imp:110}\\
79 & 79, 93, 186, 242 & (Z) & kinks of $\sh^{\pm1}F^3$; edge moves; $\Phi_f$ & \cite{PaperA}, \eqref{imp:110}\\
94 & 94, 122 & (Z) & kinks; scattering; Lemma~L; $\Phi_f$ & \cite{PaperE}, \eqref{imp:94}\\
\midrule
24 & 24, 66 & (S) & $\kappa F=F\kappa$; factor onto $46$; $\Phi_f$ & \cite{PaperA}, \eqref{imp:110}\\
189 & 189, 231 & (S) & $\kappa F=F\kappa$; factor onto $46$; $F^4=F_{66}^4$ & \cite{PaperA}, \eqref{imp:110}\\
60 & 60, 102 & (S) & affine; $\Aut=\langle\sh,F\rangle\times\langle\kappa\rangle$ & \cite{PaperC}, \eqref{imp:affine}\\
153 & 153, 195 & (S) & affine; $\Aut=\langle\sh,F\rangle\times\langle\kappa\rangle$ & \cite{PaperC}, \eqref{imp:affine}\\
\midrule
15 & 15, 85, 170, 240 & (A) & affine; $\Aut\cong\Z^2$, index $2$ over $\langle\sh,F\rangle$ & \cite{PaperC}, \eqref{imp:affine}\\
105 & 105, 150 & (A) & affine; $\Aut\cong\Z^3$, with $\mathrm{id}+\sh F$ for $150$ & \cite{PaperC}, \eqref{imp:affine}\\
\bottomrule
\end{tabular}

\medskip
\caption{Status (O), the open classes: $28$ classes, $98$ rules. For each, every automorphism of $X_{F_f}$ of radius at most $3$ is of the form $\sh^aF^b$ (Theorem~\ref{cl:thm:census}). The column $\Phi_f$ records whether $\Phi_f$ has a non-periodic point (Section~\ref{cl:sec:Z2}); for all $28$ classes $\langle\sh,F_f\rangle\cong\Z^2$ holds by Theorem~\ref{cl:thm:A}.}
\label{cl:tab:open}
\centering\footnotesize\renewcommand{\arraystretch}{0.9}
\begin{tabular}{rlcl}
\toprule
rep. & members & $\Phi_f$ & remark\\
\midrule
13 & 13, 69, 162, 176 & -- & \\
25 & 25, 67, 152, 194 & yes & \\
26 & 26, 74, 82, 88 & yes & \\
28 & 28, 56, 70, 98 & -- & \\
29 & 29, 71, 184, 226 & -- & \\
31 & 31, 87, 234, 248 & -- & \\
41 & 41, 97, 134, 148 & -- & \\
43 & 43, 113, 142, 212 & -- & \\
47 & 47, 117, 174, 244 & -- & \\
57 & 57, 99, 156, 198 & -- & permutive in the centre variable\\
58 & 58, 78, 92, 114 & yes & \\
61 & 61, 103, 188, 230 & -- & \\
77 & 77, 178 & yes & \\
91 & 91, 218 & yes & \\
107 & 107, 121, 158, 214 & -- & \\
109 & 109, 182 & -- & \\
111 & 111, 125, 190, 246 & -- & \\
123 & 123, 222 & yes & \\
141 & 141, 163, 177, 197 & -- & \\
143 & 143, 171, 213, 241 & -- & \\
151 & 151, 233 & -- & \\
155 & 155, 203, 211, 217 & yes & \\
157 & 157, 185, 199, 227 & -- & \\
159 & 159, 215, 235, 249 & -- & \\
167 & 167, 173, 181, 229 & -- & \\
175 & 175, 245 & -- & \\
183 & 183, 237 & -- & \\
191 & 191, 239, 247, 253 & -- & \\
\bottomrule
\end{tabular}
\end{table}

\subsection{The census}\label{cl:sec:census}

\begin{theorem}[computer-assisted: SAT enumeration, Section~\ref{sec:method}; one or two rules per class, the other members by Lemma~\ref{lem:classes}]\label{cl:thm:census}
Let $f$ belong to one of the $28$ classes of Table~\ref{cl:tab:open} or to the classes of $14$, $46$, $79$, $94$, $95$ and $110$, and let $R\le3$. Then every invertible block map of radius at most $R$ that commutes with $F_f$ equals $\sh^aF_f^b$ for some $(a,b)$. More precisely, the block maps of radius at most $R$ commuting with $F_f$ are
\begin{enumerate}
\item the $2R^2+2R+1$ maps $\sh^aF_f^b$ with $|a|+|b|\le R$, which are pairwise distinct;
\item the constant maps onto the constant configurations fixed by $F_f$, namely $(0,0)^{\Z}$ if $f(0,0,0)=0$ and $(1,1)^{\Z}$ if $f(1,1,1)=0$;
\item for the class of $95$ and $R=3$ only, $20$ further maps, none of them injective.
\end{enumerate}
\end{theorem}

The proof is a finite computation by the method of Section~\ref{sec:method}; the rules on which it was run and the logs are described in Appendix~\ref{app:census}. For the class of $110$ the same holds for $R=4$ (Theorem~\ref{thm:110}), and for the classes of $110$, $14$, $46$, $79$ and $94$ the first statement holds for every $R$, by \eqref{imp:110} and \eqref{imp:94}. The $20$ further maps for $95$ are identified in \cite[Proposition~6.21]{PaperE}: they are induced by block maps of the golden mean shift that conserve the number of particles and are not injective.

The census was also run for classes whose status is known. For the class of $133$ the commuting block maps of radius at most $3$ are the maps $\sh^aF^b$ with $|a|+|b|\le3$ and nothing else (no constant configuration is fixed, since $f(0,0,0)=f(1,1,1)=1$); so the automorphisms of radius $4$ of \eqref{imp:133} are the first outside $\langle\sh,F\rangle$. For the classes of $24$ and $189$ the census at radius $3$ finds exactly $50$ invertible maps, the $25$ maps $\sh^aF^b$ and their products with $\kappa$, in accordance with \eqref{imp:110}. For the classes of $63$ and $187$, of status (Z) by \eqref{imp:half}, it finds $(2R+1)^2$ invertible maps instead of $2R^2+2R+1$. The reason is that $f_{63}(a,b,c)=\neg(a\wedge b)$ does not read $c$, so $\sh^aF_{63}^b$ depends only on the cells in $[a-|b|,a]$ and can have radius smaller than $|a|+|b|$; for instance $\sh F_{63}^2$ has radius $1$. The pairs $(a,b)$ with $[a-|b|,a]\subseteq[-R,R]$ number $(2R+1)^2$, and an exact comparison of tables shows that the $49$ invertible maps found at $R=3$ are these powers; the same holds for $187$ with $[a,a+|b|]$. So the census agrees with $\Aut=\langle\sh,F\rangle$ for these classes.

Theorem~\ref{cl:thm:census} concerns radius at most $3$. A census at radius $4$ for the open classes is reported in Section~\ref{cl:sec:holdouts}; radius $5$ is out of reach of this encoding, which would need about $10^9$ clauses. The class of $95$ shows what the open status leaves room for: its censuses at radius $3$ (above) and $4$ \cite[Lemma~6.20]{PaperE} contain no automorphism outside $\langle\sh,F\rangle$, and its group contains every finite group, the smallest radius of an automorphism outside $\langle\sh,F\rangle$ being $7$ \eqref{imp:95}.

\subsection{No walls in 110R}\label{cl:sec:nowalls}

The automorphisms of status (G) that come from walls \eqref{imp:walls}--\eqref{imp:rot} all use a cell next to the permuted information whose trajectory does not depend on the context. We recall the kinds of walls of \cite{PaperC}; the fixed walls, and the periodic pairs of cells from which the other kinds are built, are due to Takesue \cite[Section~4]{Takesue1989}. Let $F$ be a reversible cellular automaton of radius $1$ on $\A^{\Z}$.
\begin{itemize}
\item A \emph{fixed wall} is a word $W$ such that $F(c)$ contains $W$ at every position at which $c$ does.
\item A \emph{time-periodic wall} is a family of distinct words $w_0,\dots,w_{P-1}\in\A^k$, indices modulo $P$, such that $c_{[0,k)}=w_s$ implies $F(c)_{[0,k)}=w_{s+1}$, for every $c$ and $s$.
\item A \emph{rotating wall} is a run of cells $c_i$ on which $F$ acts by a fixed permutation $\pi$ of $\A$ in every context, $F(c)_i=\pi(c_i)$, the run being preserved by $F$.
\item A \emph{hidden-bit wall} is given by $k\ge3$, a set $W\subseteq\A^k$, a position $j$ with $1\le j\le k-2$, and an involution $\iota\colon W\to W$ such that $\iota(w)$ differs from $w$ exactly at position $j$. Put $\Occ(c)=\{i:\ c_{[i,i+k)}\in W\}$, and for $i\in\Occ(c)$ let $\iota_ic$ be $c$ with $c_{[i,i+k)}$ replaced by $\iota(c_{[i,i+k)})$. The two basic hypotheses are \textup{(W1)}: if $0\in\Occ(c)$ then $0\in\Occ(F(c))$; and \textup{(H1)}: if $0\in\Occ(c)$ then $F(\iota_0c)=\iota_0F(c)$.
\end{itemize}
The constructions of \cite{PaperC} assume these and further hypotheses; we exclude already the weaker structures defined here. Rule $110$R has none of them, of any length.

\begin{proposition}[no walls in $110$R]\label{cl:prop:nowalls}
Let $F=F_{110}$.
\begin{enumerate}
\item For every word $w$ on an interval $[l,r]$ and every $u\in\A^{(-\infty,l)}$ there are $v,v'\in\A^{(r,\infty)}$ and $t\in\{1,2,3\}$ with $F^t(uwv)_r\ne F^t(uwv')_r$.
\item If $c$ and $c'$ differ in exactly one cell $j$, then $F^t(c)$ and $F^t(c')$ differ in some cell other than $j$ for some $t$ with $1\le t\le6$.
\item Consequently $F$ has no word whose right end cell has a trajectory independent of the context; in particular no word fixed in every context, no time-periodic wall, and no rotating wall. Moreover, for no $k\ge3$ is there a nonempty family $W\subseteq\A^k$ with an involution $\iota$ as above satisfying \textup{(W1)} and \textup{(H1)}. So $110$R has no fixed, time-periodic, rotating or hidden-bit wall of any length.
\end{enumerate}
\end{proposition}

\begin{proof}
Throughout, $f=f_{110}$, and we use three facts read off its table: $f(a,b,\cdot)$ is constant only for $(a,b)=(0,1)$, where it is $1$; $f(1,b,c)=b\oplus c$; and $f(\cdot,0,1)=f(\cdot,1,0)=1$.

(1) Write $Y$ for the diagram of $uwv$. Its values in rows $-1$ and $0$ at columns $i\le r$ do not depend on $v$. Suppose that $F^t(uwv)_r$, which is $(Y(t-1,r),Y(t,r))$, is the same for all $v$, for $t=1,2,3$. For $t=1$: $Y(1,r)=Y(-1,r)\oplus f(Y(0,r-1),Y(0,r),Y(0,r+1))$, and $Y(0,r+1)$ is the free bit $y(v_{r+1})$, so $(Y(0,r-1),Y(0,r))=(0,1)$ and $Y(1,r)=1\oplus Y(-1,r)$. For $t=2$: complementing $x(v_{r+1})=Y(-1,r+1)$ complements $Y(1,r+1)$ and changes no value in row $0$ and no $Y(1,i)$ with $i\le r$. Since $Y(2,r)=Y(0,r)\oplus f(Y(1,r-1),Y(1,r),Y(1,r+1))$, we need $(Y(1,r-1),Y(1,r))=(0,1)$. Then $Y(-1,r)=0$ and $Y(2,r)=1\oplus f(0,1,\cdot)=0$. For $t=3$: complementing $x(v_{r+2})=Y(-1,r+2)$ complements $Y(1,r+2)$ and changes no other value in rows $0$ and $1$, hence no $Y(2,i)$ with $i\le r$; and it complements $Y(2,r+1)=Y(0,r+1)\oplus f(Y(1,r),Y(1,r+1),Y(1,r+2))$, because $Y(1,r)=1$. So $Y(3,r)=Y(1,r)\oplus f(Y(2,r-1),Y(2,r),Y(2,r+1))$ is unchanged only if $(Y(2,r-1),Y(2,r))=(0,1)$, which contradicts $Y(2,r)=0$.

(2) Let $Y=Y_c$, $\Delta=Y_c\oplus Y_{c'}$ and $s_t=\Delta(t,j)$. Suppose, for a contradiction, that $\Delta(t,i)=0$ for all $i\ne j$ and $-1\le t\le T$, where $T=6$.
\begin{enumerate}
\item[(a)] $(s_{t-1},s_t)\ne(0,0)$ for $0\le t\le T$, since otherwise $F^t(c)=F^t(c')$ and $c=c'$.
\item[(b)] If $0\le\tau\le T-1$ and $s_\tau=1$, then $(Y(\tau,j-2),Y(\tau,j-1))=(0,1)$. Indeed, in the relation at $(\tau,j-1)$ the two diagrams differ only in the input $Y(\tau,j)$, and $\Delta(\tau\pm1,j-1)=0$; so $f(Y(\tau,j-2),Y(\tau,j-1),\cdot)$ is constant.
\item[(c)] $s_{\tau+1}=s_{\tau-1}\oplus s_\tau$ for $0\le\tau\le T-1$. If $s_\tau=0$, the relation at $(\tau,j)$ has the same inputs in both diagrams except $Y(\tau-1,j)$. If $s_\tau=1$, then $Y(\tau,j-1)=1$ by (b), and since $f(1,b,c)=b\oplus c$ the relation at $(\tau,j)$ gives $s_{\tau+1}\oplus s_{\tau-1}=1$.
\item[(d)] There is no $\tau$ with $2\le\tau\le T-2$ and $s_\tau=s_{\tau+1}=1$. Suppose there were. By (b), $Y(\tau,j-2)=Y(\tau+1,j-2)=0$ and $Y(\tau,j-1)=Y(\tau+1,j-1)=1$. The relation at $(\tau,j-1)$ gives $Y(\tau-1,j-1)=Y(\tau+1,j-1)\oplus f(0,1,Y(\tau,j))=0$, and the relation at $(\tau,j-2)$ gives $Y(\tau-1,j-2)=Y(\tau+1,j-2)\oplus f(Y(\tau,j-3),0,1)=1$. By (c) at $\tau$ and at $\tau-1$, $s_{\tau-1}=0$ and $s_{\tau-2}=1$, so (b) at $\tau-2\ge0$ gives $Y(\tau-2,j-2)=0$. Then the relation at $(\tau-1,j-2)$ gives $Y(\tau,j-2)=Y(\tau-2,j-2)\oplus f(Y(\tau-1,j-3),1,0)=1$, a contradiction.
\end{enumerate}
By (a), $(s_{-1},s_0)$ is $(0,1)$, $(1,0)$ or $(1,1)$, and by (c) the sequence $s_{-1},\dots,s_6$ is $0,1,1,0,1,1,0,1$, or $1,0,1,1,0,1,1,0$, or $1,1,0,1,1,0,1,1$. These have $s_\tau=s_{\tau+1}=1$ at $\tau=3$, $4$ and $2$ respectively, contrary to (d).

(3) The first statement is (1). A fixed wall, each word of a time-periodic wall, and a rotating wall all have this property: the trajectory of their right end cell is determined by the word. Next let $W\ne\emptyset$, $j$ and $\iota$ satisfy (W1) and (H1), and let $0\in\Occ(c)$. By (W1), $0\in\Occ(F^t(c))$ for all $t\ge0$, so (H1) gives by induction $F^t(\iota_0c)=\iota_0F^t(c)$ for all $t\ge0$. Since $\iota_0$ changes only the cell $j$, the configurations $\iota_0c\ne c$ differ in exactly one cell and their images under every $F^t$ differ in no other cell, contrary to (2).
\end{proof}

The proposition extends to every length the exclusion of fixed words in \eqref{imp:133} and the search for hidden-bit families in \eqref{imp:hidden}, which covers windows of length at most $6$. For words of at most $6$ cells, that the class of $110$ has no local conservation law, and hence no fixed or time-periodic wall, is part of the examination in \cite[Section~4]{Takesue1989}. It does not exclude words fixed in every context by a power $F^P$ alone, which are not one of the wall mechanisms above; a finite search (Appendix~\ref{app:stated}) finds that for $P\le6$ no word has its right end cell fixed by $F^P$ in every context. The statements for the other rules of the class of $110$ follow by reflection and complementation, with left and right exchanged in (1).

\subsection{The open classes}\label{cl:sec:holdouts}

Table~\ref{cl:tab:open} lists the classes for which none of the statuses (G), (Z), (S), (A) is known. The classes of $110$, $14$, $46$, $79$ and $94$ pass the same censuses and have $\Aut(X_{F_f})=\langle\sh,F_f\rangle$; the class of $95$ passes them as far as automorphisms are concerned and has status (G). So the absence of automorphisms outside $\langle\sh,F\rangle$ at small radius does not predict the group, and $\Aut(X_{F_f})=\langle\sh,F_f\rangle$ for the open classes is a question (Question~\ref{q:open}). We record which of the available tools bear on them.

\emph{Walls.} None of the open classes is among the $40$ classes in which Takesue exhibits walls \cite[Section~4]{Takesue1989}. No open class has a word of any length fixed in every context \eqref{imp:133}, and the search of \eqref{imp:hidden} finds hidden-bit and time-periodic families only in the classes of $27$, $59$, $127$, $179$ and $223$. We note that this search does not see families that are unions of several closures, such as the family of length $5$ for rule $126$ in \eqref{imp:rot}, and that walls fixed only by a power of $F$ with larger windows, or longer hidden-bit walls, are not excluded by it. The census would have seen the small automorphisms that walls produce in the known examples: the relabelling $D$ of $54$R has radius $2$ (Table~\ref{tab:census}), the flips of the hidden-bit walls of $59$R appear at radius $2$ (computer-assisted in \cite{PaperC}), and the flip-all maps of the time-periodic walls and of $27$R have radius $2$ and $3$ \cite{PaperC}. This is evidence that short walls of these kinds are absent from the open classes, not a proof: the wall $0330$ of $133$ has length $4$, and its first automorphism outside $\langle\sh,F\rangle$ has radius $4$, beyond the census.

\begin{sloppypar}
\emph{Exact searches.} Finite searches (Appendix~\ref{app:stated}) exclude, for the classes of Table~\ref{cl:tab:open} and the class of $95$: words of length at most $11$, $9$, $7$, $5$, $3$ that are fixed in every context by $\sh^sF^p$ with $p=1,\dots,5$ and $|s|\le p$, and invariant families of type (W1) of the same lengths, except, for $191$, the level sets of its frozen field; hidden-bit walls with words of length at most $10$; and eigenfields $\vartheta$ with $\vartheta\circ F^p=\sh^{-k}\circ\vartheta$ of radius at most $6-p$, for $p=1,\dots,5$ and $|k|\le p$, other than the frozen field of $191$. Eigenfields are the mechanism behind the non-affine commuting maps of $90$R of radius at most $3$ (computer-assisted in \cite{PaperC}). For $p=1$ and $s=0$ an invariant family of type (W1) is a local conservation law in the sense of \cite[Section~4]{Takesue1989}, and for words of length at most $6$ this part of the search is contained in the examination reported there, by which $44$ of the $88$ classes have a local conservation law on at most four cells, and the other $44$ have none on at most six. In the present classification the $44$ classes without one are the open classes other than that of $191$, the classes of status (Z) other than that of $46$, and the classes of $15$, $60$, $90$, $95$, $105$, $153$ and $165$ (Remark~\ref{cl:rem:eigen}). Takesue conjectures that four cells suffice, without proof \cite[p.~391]{Takesue1989}. What the search adds, for the classes on which it is run, is the lengths from $7$ to $11$, and the cases $p\ge2$ and $s\ne0$. The frozen field of $191$ is the local conservation law of $191$R on three cells recorded in \cite[Table~3]{HT1991}. The same searches recover additive conserved quantities for $26$, $77$, $91$, $123$, $175$ and $191$ among the open classes, and for $46$, $94$ and $95$. For densities on at most three adjacent cells these are the quantities of \cite[Tables~3 and~4]{HT1991} (see also \cite{Takesue1989,Takesue1995,Takesue1997}); the quantity of $175$ has a density on four cells, which is beyond the range of \cite{HT1991}; its existence is recorded in \cite[Table~1]{Takesue1997}, which lists one quantity for $175$R with density on four cells and none on fewer, without writing out the density. None of them is known to give an automorphism in an open class. For $95$ they count the particles of each family on each sublattice \cite{PaperE}, and for $46$ the frozen fields are the bond bits used in \cite{PaperA}, that is, the local conservation law of $46$R of \cite[(4.4)--(4.6)]{Takesue1989}.
\end{sloppypar}

\emph{Radius $4$.} A census at radius $4$ (computer-assisted; Appendix~\ref{app:census}), with the method of Section~\ref{sec:method}, covers one rule of each open class and of the classes of $14$, $24$, $46$, $63$, $79$, $94$, $110$, $126$, $133$, $187$ and $189$, $219$. For the open classes and for $14$, $46$, $79$, $94$ and $110$ the commuting block maps of radius at most $4$ are the $41$ maps $\sh^aF^b$ and the constants. For $24$ and $189$ the further solutions are the maps $\kappa\sh^aF^b$, and for $63$ and $187$ they are powers $\sh^aF^b$ of radius smaller than $|a|+|b|$, as an exact comparison of tables confirms. For $133$ there are $27$ further solutions, which are not classified here; automorphisms of radius $4$ outside $\langle\sh,F\rangle$ exist by \eqref{imp:133}. For $126$ and $219$ the enumeration reached its cap. Injectivity is not tested at this radius. For the class of $95$ the census at radius $4$ is \cite[Lemma~6.20, Proposition~6.21(a)]{PaperE}, where it is carried out with two solvers: there are $262$ commuting maps, $220$ of them neither of the form $\sh^aF^b$ nor constant, and none of these is injective.

\emph{The form of the rule.} No open class is permutive in an outer variable: the non-affine outer-permutive rules are the four classes of \eqref{imp:perm}, and the affine ones are the classes of $15$, $105$, $60$, $153$, $90$ and $165$. The rigidity theorem of \cite{PaperB} for automata of any radius applies to rules that are permutive in an outer variable, so it would apply to an open class only through a recoding that makes it outer-permutive, and we know none. The class of $57$ is the only open class that is permutive in the centre variable; the centre-permutive rules are the classes of $51$, $54$, $57$, $60$, $105$, $147$ and $153$. No open class contains a rule that ignores an outer variable: the $28$ such rules lie in classes of status (G), (Z), (S) or (A), and \eqref{imp:half} covers the non-affine ones without fixed walls. The rules that do not read their middle argument are the $16$ rules of the classes of $0$, $5$, $10$, $15$, $90$, $95$, $165$, $175$ and $255$ \cite[p.~386]{Takesue1989}; the class of $175$ is the only open one, and for it see \cite[Remark~6.22]{PaperE}.

\emph{Light-speed fixed points.} The proofs of status (Z) in \cite{PaperA,PaperE} begin with a homomorphism $\Aut(X_{F_f})\to\Z^2$ that extends $\sh^aF^b\mapsto(a,b)$, read off from the action of an automorphism on kinks among the fixed points of $\sh^{\pm1}F_f$, that is, on non-periodic points of $\Phi_f$. For $7$ of the open classes $\Phi_f$ has non-periodic points (Table~\ref{cl:tab:open}): for the class of $25$ they come from a path between two cycles, as for $110$, and for the other six $\Phi_f$ has positive entropy, as for $94$ and $95$. For the remaining $21$ open classes $\Phi_f$ is finite, and for $9$ of them ($141$, $143$, $151$, $157$, $159$, $167$, $175$, $183$, $191$) it is empty: there $f(0,0,0)=f(1,1,1)=1$ and no constant configuration is fixed by $F$. For these $21$ classes no such homomorphism is known.

\emph{Blocking words.} In the wall constructions of \cite{PaperC} the automorphisms outside $\langle\sh,F\rangle$ come from a blocking word (Section~\ref{sec:blocking}) with an autonomous degree of freedom next to it. The absence of blocking words does not imply that $\Aut(X_{F_f})$ is virtually $\Z^2$: not for affine rules, by \eqref{imp:90}, and not for rules that are not affine, by \eqref{imp:95}. It is not known which open classes have blocking words.

\emph{Window graphs and eigenfields.} By Lemma~\ref{cl:lem:eigen}, the condition $(\Lambda^{\pm})$ of Section~\ref{sec:blocking} says that $F$ has no eigenfield of speed $\pm1$, which the eigen-factors of $90$R violate \cite{PaperC}, and the zero-speed condition says that $F$ has no frozen field. For $95$R the graphs $W^{(k,p)}_M$ are strongly connected for all $(k,p)$ and all $M$ \eqref{imp:95}.

\begin{remark}[window graphs of the classes; computer-assisted]\label{cl:rem:eigen}
A computation of the graphs $W^{(\pm1,1)}_M$ and $W^{(0,1)}_M$ for $M\le6$ (Appendix~\ref{app:stated}) gives the following. Every class of status (G) except those of $90$, $165$ and $95$ fails the zero-speed condition at some $M\le3$; $90$ and $165$ satisfy it for $M\le6$ and fail $(\Lambda^{\pm})$. The classes of status (Z) other than that of $46$, and $27$ of the $28$ open classes, satisfy both conditions for $M\le6$; the exceptions $46$ and $191$ fail the zero-speed condition and have frozen fields. In status (S), the classes of $24$ and $189$ fail the zero-speed condition, at $M=1$ and $M=2$, and the affine classes of $60$ and $153$ satisfy both conditions for $M\le6$, as do the classes of $15$ and $105$ of status (A). The class of $95$ satisfies both conditions for every $M$, by \eqref{imp:95}. For the zero-speed graphs with $M\le5$, that is, for words of at most $6$ cells, these statements reproduce the examination of local conservation laws in \cite[Section~4]{Takesue1989}; the graphs of speed $\pm1$ are not considered there. These computations are not used elsewhere.
\end{remark}

\subsection{Reach of the edge-move method}\label{cl:sec:reach}

The proof of \eqref{imp:110} in \cite{PaperA} works at the right edge of configurations that agree with a spatially periodic background on a right half-line. We record where this method stops on the open classes. Let $f$ be an elementary rule and $F=F_f$. A configuration $B$ (in the applications, a spatially periodic one) has a \emph{stuck right edge at $r$} if for every $t\in\Z$ and every configuration $c$ that agrees with $F^tB$ on $(r,\infty)$ and differs from it at $r$, the configurations $F(c)$ and $F^{-1}(c)$ agree with $F^{t+1}B$ and $F^{t-1}B$ on $(r,\infty)$ and differ from them at $r$. Stuck left edges are defined in the same way with $(-\infty,r)$, and the reflection $\mathrm{rev}$ exchanges the two. For $110$, $42$ and $46$ the constant $(0,0)^{\Z}$ has a stuck right edge at every $r$, and for $93$ a background of spatial period $2$ has one \cite{PaperA}. Let
\[
N_f=\{(b,c)\in\{0,1\}^2:\ f(0,b,c)=f(1,b,c)\}
\]
be the set of contexts in which $f$ ignores its left input, and for $(b,c)\in N_f$ let $g_f(b,c)$ be the common value.

\begin{lemma}\label{cl:lem:stuckcols}
Let $B$ have a stuck right edge at $r$, and put $u_t=Y_B(t,r+1)$, $v_t=Y_B(t,r+2)$, $w_t=Y_B(t,r+3)$, $z_t=Y_B(t,r+4)$. Then for every $t\in\Z$:
\begin{enumerate}
\item $(u_t,v_t)\in N_f$;
\item $u_{t+1}\oplus u_{t-1}=g_f(u_t,v_t)$;
\item $v_{t+1}\oplus v_{t-1}=f(u_t,v_t,w_t)$ and $w_{t+1}\oplus w_{t-1}=f(v_t,w_t,z_t)$.
\end{enumerate}
\end{lemma}

\begin{proof}
Fix $t$, let $d=F^tB$, and let $c$ be $d$ with the bit $y(d_r)$ complemented. Then $c$ agrees with $d$ on $(r,\infty)$ and differs from it at $r$, so $F(c)_{r+1}=F(d)_{r+1}$. By \eqref{eq:F} the second bits of these cells are $x(d_{r+1})\oplus f(\neg y(d_r),y(d_{r+1}),y(d_{r+2}))$ and $x(d_{r+1})\oplus f(y(d_r),y(d_{r+1}),y(d_{r+2}))$, and the $y$-row of $d$ is row $t$ of $Y_B$; this gives (1). Parts (2) and (3) are the spacetime relation \eqref{eq:rel} of $Y_B$ at $(t,r+1)$, $(t,r+2)$ and $(t,r+3)$, with (1) used in (2).
\end{proof}

\begin{proposition}\label{cl:prop:nostuck}
Let $f$ be a rule in one of the $12$ classes of $77$, $91$, $94$, $95$, $107$, $109$, $123$, $141$, $151$, $167$, $175$ and $183$. Then no configuration has a stuck right edge or a stuck left edge for $F_f$; in particular no spatially periodic one has.
\end{proposition}

\begin{proof}
The complement $\kappa$ commutes with $\sh$ and conjugates $F_f$ to $F_{\kappa f}$, so it maps a configuration with a stuck right edge for $F_f$ to one for $F_{\kappa f}$; and $\mathrm{rev}$ maps stuck left edges for $F_f$ to stuck right edges for $F_{f^r}$, which lies in the same class. So it suffices to exclude stuck right edges for one rule of each pair $\{f,\kappa f\}$ in these classes, namely for
\[
91,\ 94,\ 95,\ 107,\ 121,\ 109,\ 123,\ 151,\ 167,\ 173,\ 175,\ 183,\ 77,\ 141,\ 163.
\]
Suppose that $B$ has a stuck right edge at $r$, and use the notation of Lemma~\ref{cl:lem:stuckcols}. The values of $f$ used below are read off its table.

(a) For the first twelve rules, $g_f\equiv1$ on $N_f$, so by (2) $u_{t+1}\ne u_{t-1}$ for every $t$, and $u$ takes both values. For $91$, $94$, $107$, $121$, $109$, $123$, $151$, $167$, $173$, $183$ the set $N_f$ is $\{00\}$, $\{10\}$, $\{01\}$, $\{00\}$, $\{10\}$, $\{00,01\}$, $\{00\}$, $\{01\}$, $\{11\}$, $\{00,01\}$, whose elements share their first coordinate, so $u$ is constant by (1), a contradiction. For $95$, $N_f=\{00,10\}$, and for $175$, $N_f=\{01,11\}$; in both cases $v$ is constant by (1), so $v_{t+1}\oplus v_{t-1}=0$, and at a time $t$ with $u_t=0$ part (3) requires $f(0,v_t,w_t)=0$, while $f_{95}(0,0,\cdot)=1$ and $f_{175}(0,1,\cdot)=1$.

(b) For $77$ and $163$, $N_f=\{01,10\}$, so $v_t=\neg u_t$ and (2), (3) give $f(u_t,v_t,w_t)=g_f(u_t,v_t)$. At a time with $u_t=0$ this reads $f(0,1,w_t)=g_f(0,1)$, while $f_{77}(0,1,\cdot)=1\ne0=g_{77}(0,1)$ and $f_{163}(0,1,\cdot)=0\ne1=g_{163}(0,1)$. So $(u_t,v_t)=(1,0)$ for all $t$, and (2) requires $g_f(1,0)=0$. For $77$, $g_{77}(1,0)=1$. For $163$, $g_{163}(1,0)=0$, and (3) gives $f_{163}(1,0,w_t)=0$ for all $t$, so $w\equiv0$ since $f_{163}(1,0,1)=1$; then the second relation in (3) reads $0=f_{163}(0,0,z_t)$, while $f_{163}(0,0,\cdot)=1$.

(c) For $141$, $N_f=\{01,11\}$, so $v\equiv1$ and $f_{141}(u_t,1,w_t)=0$ by (3). Since $f_{141}(0,1,\cdot)=1$, $u\equiv1$, and (2) requires $g_{141}(1,1)=0$, while $g_{141}(1,1)=1$.
\end{proof}

In these $12$ classes, of which all but those of $94$ and $95$ are open, the construction of \cite{PaperA} cannot start: there is no edge at which to glue a kink. They are not the only classes without stuck edges: for the affine classes of $90$, $105$ and $165$ the rule is permutive in both outer variables, $N_f$ is empty, and part (1) of Lemma~\ref{cl:lem:stuckcols} excludes stuck edges at once. For the remaining open classes the obstruction is located by computation.

\begin{proposition}[computer-assisted]\label{cl:prop:reach}
Let $f$ be a rule of one of the $11$ classes of $28$, $29$, $41$, $43$, $47$, $57$, $61$, $111$, $143$, $155$ and $157$, let $B$ be a spatially periodic configuration of spatial period at most $6$ with a stuck right edge, and let $1\le b\le4$ and $0<|a|\le b$. If $B\in\mathrm{Fix}(\sh^aF_f^b)$, then $B$ is the only point of $\mathrm{Fix}(\sh^aF_f^b)$ that agrees with $B$ on a left half-line.
\end{proposition}

The proof is an exact finite search (Appendix~\ref{app:reach}): in the graph of the windows of $\mathrm{Fix}(\sh^aF_f^b)$, every window of $B$ has exactly one outgoing edge. In particular, for these rules, periods and speeds there is no kink with such a background on its left, and nothing to glue next to the edge.

\begin{remark}[sampled computations]\label{cl:rem:reach}
The following computations use sampled left contexts (Appendix~\ref{app:reach}); they are observations, not proofs. In \cite{PaperA}, a kink glued next to a stuck edge leaves, after a bounded time, a configuration with the same edge that is changed only near the edge; this defines a \emph{move}, a partial map on the states of the last $k$ cells before the edge, and the proof needs the orbits of the moves to reach the states from which the next column to the left can be read.
\begin{enumerate}
\item For the class of $159$ there are kinks of speed $1/2$ with a background with a stuck edge on their left, but no kink glued next to the edge was found to give a move.
\item For the classes of $13$, $25$, $26$, $58$ and $191$ there are backgrounds with a stuck edge, kinks and moves, but the orbits of the moves on the last $k\le4$ cells do not reach the states from which the next column is read.
\item For the class of $31$ there are backgrounds, kinks of speed $1/3$ and moves, and for every $k\le6$ the moves on the last $k$ cells have two orbits for one value of the cell beyond the edge and a single orbit for each other value, as for $110$. A proof that a block map commuting with $F_{31}$ and fixing the kinks commutes with the moves is missing.
\end{enumerate}
With Propositions~\ref{cl:prop:nostuck} and~\ref{cl:prop:reach} this accounts for the $28$ open classes: $10$, $11$, $1$, $5$ and $1$.
\end{remark}

\section{Discussion and open questions}\label{sec:discussion}

\subsection{The first-order rule is no guide}
The first-order rules $54$ and $110$ are the elementary rules associated with class~4 behaviour \cite[pp.~693--697]{Wolfram2002}; their second-order versions have automorphism groups of opposite kinds, and $110$R and $30$R, whose first-order rules behave differently, have the same group $\Z^2$ and the same commuting block maps through radius $4$. What distinguishes $54$R is a conserved local structure, the wall $33$, with an autonomous interior. The same mechanism, in the forms of fixed, hidden-bit, time-periodic and rotating walls, accounts for $39$ of the $43$ classes of status (G); the walls of these classes go back to Takesue \cite[Section~4]{Takesue1989}, and their use for automorphisms is that of \cite{PaperC}. Three of the other four, $90$, $165$ and $255$, are affine \cite{PaperC}. The fourth is the class of $95$, which is not affine: there the room comes from the particle-conserving automorphisms of a gas of light-speed particles, and $F_{95}$ has no blocking word and no eigenfield of any speed \eqref{imp:95}. The particles of the first-order rule $54$ are well studied \cite{BNR1991,MACC2012}; they play no role here. The reversible automaton called Rule~54 in \cite{BBGP1993,BKP2021} is neither rule $54$ nor $54$R: it is the system on one sublattice of a spacetime diagram of the class of $95$ \cite{PaperE}.

\subsection{Rigidity}
On the open classes the edge-move argument of \cite{PaperA} does not get started in $21$ classes, for lack of a stuck edge or of a kink next to a background with one, and stops later in the computations for $6$ more (Section~\ref{cl:sec:reach}). For $91$ Lemma~L holds at every radius, but no normalisation of the automorphisms is known; for $155$ even the proof of Lemma~L fails \eqref{imp:L}. For the $21$ open classes with finite $\Phi_f$ there are no kinks of light speed from which to read a homomorphism to $\Z^2$.

Proposition~\ref{cl:prop:nowalls} shows that $110$R has no fixed, time-periodic, rotating or hidden-bit wall of any length; the rigidity of $110$R \cite{PaperA} is the stronger statement that it has no automorphism outside $\langle\sh,F\rangle$ of any kind. Conjugacy of reversible automata is undecidable \cite{Salo2020}; we do not know whether the property $\Aut(X_F)=\langle\sh,F\rangle$ of a reversible automaton $F$ is decidable.

\subsection{Open questions}
A word fixed in every context is blocking (Section~\ref{sec:blocking}), and the walls behind the status (G) entries other than $90$, $95$, $165$ and $255$ all give blocking words \cite{PaperC}; for $255$ every word is blocking. Among the $43$ classes of status (G), those of $90$, $165$ and $95$ are the only ones without a blocking word \eqref{imp:90}, \eqref{imp:95}. Question~1.9 of \cite{PaperC} asks whether, for an elementary rule $f$ that is not affine and does not belong to the class of $95$, the absence of a blocking word for $F_f$ implies that $\Aut(X_{F_f})$ is virtually $\Z^2$. The classification answers it for the classes of status (G), (Z), (S) and (A): the classes of status (A) are affine, the classes of status (G) that are not affine, other than that of $95$, have blocking words, and in the classes of status (Z) and (S) the group is virtually $\Z^2$. So that question concerns the $28$ open classes.

\begin{question}\label{q:open}
Is $\Aut(X_{F_f})=\langle\sh,F_f\rangle$ for each of the $28$ open classes?
\end{question}

The evidence is weak: the class of $95$ passes the same censuses. The natural candidate for a second exception is the class of $175$, the only open class whose rule does not read its middle argument, so that the two sublattices of a diagram evolve independently, as for $95$ \cite[Remark~6.22]{PaperE}.

\begin{enumerate}
\item Decide the $28$ open classes. Possible routes: for the class of $25$, the remaining open class whose light-speed fixed points include kinks as for $110$, and for the class of $31$, where kinks of $\sh^{\pm1}F^3$ and moves behave as for $110$ in the computations of Remark~\ref{cl:rem:reach}, an argument as in \cite{PaperA}; for $91$, a normalisation of the automorphisms as for $94$, and for $155$, a proof of Lemma~L (see the questions in \cite{PaperE}). For the $21$ open classes with finite $\Phi_f$ a homomorphism $\Aut(X_{F_f})\to\Z^2$ extending $\sh^aF^b\mapsto(a,b)$ is missing.
\item Extend the census beyond radius $4$ for the open classes.
\item The following are questions of \cite{PaperC}. Determine $\Aut(X_{F_f})$ for a class of status (G) with walls. For the affine classes of statuses (S) and (A), is every block map commuting with $F_f$ affine? Every automorphism is \eqref{imp:affine}, and every commuting block map of radius at most $3$ is (computer-assisted in \cite{PaperC}). For $54$R, is $\Aut(X_{F_{54}})$ generated by $\sh$, $F$ and the wall relabellings and wall-trading involutions of the kind in \cite{PaperC}, of all radii? For $90$R, is it generated by the affine automorphisms and the marker maps of \cite{PaperC}?
\item For the class of $95$, describe the group $\mathcal N$ of \eqref{imp:95} and the block maps that commute with $F_{95}$ and are not invertible; see the questions in \cite{PaperE}.
\end{enumerate}


\appendix
\section{Computations}\label{app:comp}

This appendix describes the finite computations. Every result marked as computer-assisted in the text is the output of a terminating exact search over a finite set; where random sampling is used, this is said, and such runs serve only as checks. The directory \texttt{code/} contains the programs, with instructions for running them and their expected outputs, and the logs of the long runs; the program files are named below next to each computation. Solution tables and other data files of more than $5$\,MB are not included; the programs that produce them are.

\subsection{The censuses}\label{app:census}

The censuses of Theorems~\ref{thm:110} and~\ref{cl:thm:census}, and the control rows of Table~\ref{tab:census} for $54$R, $90$R and $150$R, use the encoding, the enumeration, the re-verification and the invertibility test of Section~\ref{sec:method}. The programs and their logs are in \texttt{code/census}; the raw solution tables are not included, and \texttt{enum\_r4.py} produces those of radius $4$ again.

\emph{Coverage of Theorem~\ref{cl:thm:census}.} The enumeration (\texttt{code/census/centralizer.py}, with a cap of $5000$ solutions) is run at $R=1,2,3$ on $72$ rules, which represent $55$ of the $57$ classes without a fixed word of length at most $3$: all except the classes of $223$ and $255$, whose status is decided otherwise. Each class is represented by its smallest rule, and the classes of $13$, $14$, $15$, $25$, $26$, $28$, $29$, $30$, $41$, $43$, $57$, $58$, $62$, $75$, $77$, $94$ and $105$ also by the rules $162$, $42$, $170$, $152$, $74$, $56$, $184$, $106$, $134$, $142$, $156$, $78$, $110$, $154$, $178$, $122$ and $150$; this suffices by Lemma~\ref{lem:classes}. The summary lines of all runs are in \texttt{code/census/logs/census\_R123\_summary.txt}. For all $28$ classes of status (O), and for the classes of $14$, $46$, $63$, $79$, $94$, $95$, $110$, $133$ and $187$, the enumeration finishes below the cap at every $R\le3$, so the enumeration is complete there. For $133$ the solutions at $R\le3$ are the maps $\sh^aF^b$ only. For the classes of $30$, $45$, $75$ and $135$ the solutions at $R\le3$ are the maps $\sh^aF^b$ and the constants; for $24$ and $189$ there are in addition $5$, $13$, $25$ invertible maps at $R=1,2,3$, the maps $\kappa\sh^aF^b$; for $63$ and $187$ there are in addition $4$, $12$, $24$ invertible maps, the radius-deficient powers; for $27$ the solutions at $R\le2$ are the $13$ maps $\sh^aF^b$ and one constant map. The cap is reached only for classes whose status is decided otherwise: those of $59$, $127$, $179$, $90$, $15$ (run as $15$ and $170$), $60$, $153$, $165$, $126$ and $219$. For $27$ ($1742$ solutions at $R=3$) and for $105$ and $150$ ($4096$ solutions each at $R=3$) the enumeration finishes below the cap, and the solutions at $R=3$ are not classified; these classes have status (G) and (A), and Theorem~\ref{cl:thm:census} does not concern them. Tables~\ref{cl:tab:finite}--\ref{cl:tab:open} are assembled from the conjugacy classes, the fixed words of length at most $3$, the tests for affine and outer-permutive rules and the search over the $23$ nontrivial cell permutations (\texttt{code/census/classes88.py}), the hidden-bit and time-periodic walls of \cite{PaperC}, and the census output.

\emph{Radius $4$.} The census at radius $4$ of Section~\ref{cl:sec:holdouts} uses the same encoding with MiniSat~2.2 and a cap of $1000$ solutions (\texttt{enum\_r4.py}; summary lines in \texttt{logs/census\_R4\_summary.txt}), on one rule of each of the $28$ open classes and on the rules $14$, $24$, $46$, $62$, $63$, $79$, $94$, $126$, $133$, $187$, $189$ and $219$. Every solution is rechecked on all $4^{11}$ windows and matched against the $41$ maps $\sh^aF^b$ and the constants; injectivity is not tested at this radius. For the open classes and for $14$, $46$, $62$, $79$ and $94$ there is no further solution; for $133$ there are $27$, and for $126$ and $219$ the cap is reached. The $41$ further solutions for $24$ and $189$ are the maps $\kappa\sh^aF^b$, and the $40$ for $63$ and $187$ are compared exactly with the tables of the powers $\sh^aF^b$ whose neighbourhood lies in $[-4,4]$ (\texttt{identify\_r4.py}). The log also contains a run for rule $95$, which reproduces the count of $262$ solutions of \cite[Lemma~6.20]{PaperE}.

\emph{Cross-check.} The program \texttt{crosscheck.py} repeats the enumerations with other solvers, compares the solution sets with one another as sets of tables and with the logs of the enumerations (numbers of clauses and of solutions, the counts of the maps $\sh^aF^b$, of the constants and of the others, and the set itself whenever the log determines it), and rechecks every table on all windows. At $R\le3$ it covers the $72$, $67$ and $61$ instances at $R=1,2,3$ that finish below the cap, with Glucose~4, MiniSat~2.2 and CaDiCaL~1.5.3 (log \texttt{logs/crosscheck\_R123.txt}); at $R=4$ the rules $110$ and $30$ with Glucose~4 and CaDiCaL~1.5.3 (\texttt{logs/crosscheck\_R4\_110\_30.txt}) and the other $39$ rules that finish below the cap with Glucose~4 (\texttt{logs/crosscheck\_R4.txt}). There is no disagreement. With the option \texttt{--rup} the formula consisting of the encoding and one blocking clause for each solution is solved by Glucose~4 with proof logging, and the proof is checked lemma by lemma by reverse unit propagation, by a checker of about $70$ lines in the same file; this was done for all instances at $R\le3$ except those of the rules $27$, $105$ and $150$ at $R=3$, which have $1742$, $4096$ and $4096$ solutions. The cross-check was run on a machine with 94\,GB of memory, on which CaDiCaL solves the instances of radius $4$.

\emph{Rule $90$ at radius $2$.} For the control row of Table~\ref{tab:census}, the invertibility of each of the $2^{17}$ commuting maps was tested with the pair-graph test, in $5{,}851$ seconds: $26$ are invertible, all affine (\texttt{linear\_90.py}, \texttt{analyze\_90.py}; log \texttt{logs/analyze\_90\_R2.txt}).

\subsection{Computer-assisted statements}\label{app:stated}

\emph{Walls and searches} (Section~\ref{cl:sec:holdouts} and Proposition~\ref{cl:prop:nowalls}; programs and logs in \texttt{code/searches}). The searches are run on one rule of each of the $28$ open classes and on the rules $14$, $46$, $62$, $63$, $79$, $94$, $95$, $110$, $133$ and $187$. \emph{Fixed and time-periodic words, families of type (W1)} (\texttt{w1.c}): for $(K,p)=(11,1),(9,2),(7,3),(5,4),(3,5)$ and all $|s|\le p$, the graph on $\A^K$ with an edge from each word to each of its images under $\sh^sF^p$ in all contexts is built, and its strongly connected components that no edge leaves are computed; a proper invariant family, in particular a word fixed by $\sh^sF^p$ in every context, exists if and only if there is such a component other than $\A^K$. There is none, except at $s=0$ for $46$ and $191$ (level sets of their frozen fields) and for $133$ (its wall). For $p=1$, $s=0$ and $K\le6$ this is contained in the examination of local conservation laws in \cite[Section~4]{Takesue1989}; the frozen field of $46$ is the local conservation law of \cite[(4.4)--(4.6)]{Takesue1989}, that of $191$ is recorded in \cite[Table~3]{HT1991}, and the wall of $133$ is that of \cite[p.~391]{Takesue1989}. \emph{Hidden-bit walls} (\texttt{hb.c}): for $K=10$ and every position $1\le j\le8$, the largest set of pairs of words that differ exactly in the cell $j$ and whose images under $F$ in every context are again such pairs is computed; it is empty except for $133$. \emph{Eigenfields} (\texttt{factor.c}): for $(r,p)=(5,1),(4,2),(3,3),(2,4),(1,5)$ and all $|k|\le p$, the finest partition of $\A^{2r+1}$ compatible with $\vartheta\circ F^p=\sh^{-k}\circ\vartheta$ is computed by union--find, and likewise for binary $\vartheta$ with $\vartheta\circ F^p=\neg\,\sh^{-k}\circ\vartheta$; a nonconstant solution exists only at $k=0$, for $46$, $133$ and $191$. \emph{Conserved quantities} (\texttt{conserved.c}, \texttt{cq.py}): the additive conserved quantities of $F_f$ with a density $h$ on $r$ adjacent cells are the solutions of the linear system $h(F(w))-h(w_1\cdots w_r)=g(w_1\cdots w_{r+1})-g(w_0\cdots w_r)$, $w\in\A^{r+2}$, in the unknowns $h$ and $g$, where $F(w)$ is the image of the $r$ inner cells; it is solved for densities on $r\le5$ adjacent cells, modulo the prime $32749$ (which bounds the dimension over $\mathbb Q$ from above) and for $r\le3$ over $\mathbb Q$; nontrivial densities exist for $26$, $46$, $77$, $91$, $94$, $95$, $123$ (from $r=2$), $191$ (from $r=3$) and $133$, $175$ (from $r=4$), and for no other rule of the list. For densities on at most three adjacent cells this agrees with \cite[Table~3]{HT1991} and \cite[Table~5]{Takesue1995}, which do not go further; it is consistent with \cite{Takesue1997}. For $110$R (\texttt{nowalls110.py}), the right end cell of a word is fixed by $F^P$ in every context for no word and no $P\le6$.

\emph{Reach of the edge-move method} (Proposition~\ref{cl:prop:reach}, Remark~\ref{cl:rem:reach}).\label{app:reach} For the proposition, for each of the $44$ rules of the $11$ classes, all spatially periodic configurations of spatial period at most $6$ are enumerated with their $F$-orbits, and those are kept that satisfy condition (1) of Lemma~\ref{cl:lem:stuckcols} at some pair of adjacent columns in every phase of the orbit; this includes every background with a stuck right edge. For $1\le b\le4$ and $0<|a|\le b$, the window graph of $\mathrm{Fix}(\sh^aF^b)$ has as vertices the words of length $2b$ and as edges the words $w$ of length $2b+1$ whose image under $F^b$, a single cell, equals the cell $w_{b+a}$; it is pruned to the vertices on bi-infinite paths. For every kept background $B$ in $\mathrm{Fix}(\sh^aF^b)$, the program checks that every window of $B$ has exactly one outgoing edge (\texttt{code/reach/reach1\_unique.py}); a second program computes the strongly connected components and checks that no path from a window of $B$ reaches a recurrent component other than that of $B$ (\texttt{reach1.py}, which needs SciPy). Both find the known kinks for the classes of $110$, $14$ and $79$, and kinks for $25$, $31$ and $159$. For the remark (\texttt{code/reach/moves.py}; log \texttt{moves\_log.txt}), kinks with a stuck-edge background on their left are glued next to the edge of configurations that agree with the background to the right of the edge, and the glued configuration is followed in time; a move on the last $k$ cells and the exterior cell is recorded when the result is a configuration of the same kind changed only on these cells. The moves are tabulated on all $4^{k+1}$ states of these cells, with the cells further left drawn at random in $4$ independent contexts, a state being used only if all four give the same result, and their orbits are computed by union--find. The runs are: rule $87$ (class of $31$), background $(0,0)$, kinks of speed $1/3$, $k\le6$; rule $162$ (class of $13$), the background with rows $x=(00)^\infty$, $y=(01)^\infty$, speed $1/3$, $k\le4$; rule $25$, the background with rows $x=(01)^\infty$, $y=(10)^\infty$, speed $1$, $k\le4$; rule $74$ (class of $26$), background $(0,0)$, speed $1$, $k\le4$; rule $114$ (class of $58$), background $(1,1)$, speed $1$, $k\le4$; rule $191$, background $(0,0)$, speed $1/2$, $k\le4$; and rule $159$, background $(0,0)$, speed $1/2$, where no kink gives a move. As controls, the program reproduces the orbits of the moves described in \cite{PaperA} for $110$ ($k\le7$), $42$ and $46$ ($k\le8$).

\emph{Window graphs} (Remark~\ref{cl:rem:eigen}). The graphs $W^{(0,1)}_M$, $W^{(1,1)}_M$ and $W^{(-1,1)}_M$ are built from the cell rule for $M\le6$, for the smallest rule of each of the $88$ classes, and their strongly connected components are computed (\texttt{code/windows/window\_graphs.py}; output \texttt{window\_graphs\_M0-6.txt}).

\subsection{Confirmations of results proved by hand}\label{app:confirm}

The following computation checks a result that is proved by hand. It is not used in a proof.

\emph{Rule $110$R, no walls} (Proposition~\ref{cl:prop:nowalls}). For (1), over all $4^4$ values of $c_{[r-3,r]}$, the first $t$ at which the right context matters is $1$, $2$ and $3$ for $192$, $48$ and $16$ of them. For (2), over all windows and all three changes of a cell, $147{,}456$ pairs have a difference that stays in its column through $t=4$ and none through $t=5$, so $t\le5$ suffices (\texttt{code/searches/nowalls110.py}, exhaustive).

\subsection{Running times and reproducibility}\label{app:repro}

The censuses were run on one core of a machine with 32\,GB of memory under Linux, using Python~3.12 or later, NumPy, SciPy and PySAT~1.9 with CaDiCaL~1.5.3, MiniSat~2.2, Glucose~4 and MiniCard. Table~\ref{tab:times} lists instance sizes and times. The solution sets for rules $110$, $30$, $54$ and $90$ were each produced in two runs, with identical results. The radius-$4$ instances were solved with MiniSat~2.2, which needed about 5\,GB; CaDiCaL exceeded 12\,GB on the same instance. The 54R radius-$2$ invertible census used MiniCard's native cardinality constraints for the balance condition.

\begin{table}[h]
\caption{Instance sizes and running times (encode / enumerate / classify, in seconds).}
\label{tab:times}
\centering
\begin{tabular}{lrrrl}
\toprule
instance & variables & clauses & solutions & time (s)\\
\midrule
110R, 30R, $R=1$ & $128$ & $1.8\times10^4$ & 7 & $<1$\\
110R, 30R, $R=2$ & $2{,}048$ & $2.9\times10^5$ & 15 & $0.5$ / $0.3$ / $1$\\
110R, 30R, $R=3$ & $32{,}768$ & $4.7\times10^6$ & 27 & $7$ / $15$--$22$ / $33$\\
110R, 30R, $R=4$ (MiniSat) & $524{,}288$ & $7.5\times10^7$ & 43 & $130$--$160$ / $70$ / $47$\\
54R, $R=2$, balanced & $2{,}048+4{,}096$ & $2.9\times10^5$+card. & 23 & $3$ / $11$ / --\\
90R, $R=3$, elimination & $4{,}096$ classes & $2.6\times10^5$ eqs & $2^{65}$ & $4$\\
\bottomrule
\end{tabular}
\end{table}


\medskip
\noindent\textbf{Use of AI tools.} Large language models (Anthropic's Claude models) were used extensively in this work: to propose and develop proofs, to design, write and run the computer searches and their independent re-implementations, to check proofs and computations, and to draft the manuscript. The models are not authors. The author has reviewed the content and takes full responsibility for it. All code and computational output accompany the paper (Appendix~\ref{app:comp}).

\begin{thebibliography}{99}

\bibitem{Glucose} G.~Audemard and L.~Simon, Predicting learnt clauses quality in modern SAT solvers, in \emph{Proc. 21st Int. Joint Conf. on Artificial Intelligence (IJCAI 2009)}, 2009, 399--404.

\bibitem{BRY2018} M.~Baake, J.~A.~G.~Roberts and R.~Yassawi, Reversing and extended symmetries of shift spaces, \emph{Discrete Contin. Dyn. Syst.} \textbf{38} (2018), 835--866.

\bibitem{CaDiCaL} A.~Biere, K.~Fazekas, M.~Fleury and M.~Heisinger, CaDiCaL, Kissat, Paracooba, Plingeling and Treengeling entering the SAT Competition 2020, in \emph{Proc. SAT Competition 2020}, Univ. Helsinki, 2020, 51--53.

\bibitem{BT2000} F.~Blanchard and P.~Tisseur, Some properties of cellular automata with equicontinuity points, \emph{Ann. Inst. H. Poincar\'e Probab. Statist.} \textbf{36} (2000), 569--582.

\bibitem{BBGP1993} A.~Bobenko, M.~Bordemann, C.~Gunn and U.~Pinkall, On two integrable cellular automata, \emph{Comm. Math. Phys.} \textbf{158} (1993), 127--134.

\bibitem{BNR1991} N.~Boccara, J.~Nasser and M.~Roger, Particlelike structures and their interactions in spatiotemporal patterns generated by one-dimensional deterministic cellular-automaton rules, \emph{Phys. Rev. A} \textbf{44} (1991), 866--875.

\bibitem{Boykett2004} T.~Boykett, Efficient exhaustive listings of reversible one dimensional cellular automata, \emph{Theoret. Comput. Sci.} \textbf{325} (2004), 215--247.

\bibitem{BLR1988} M.~Boyle, D.~Lind and D.~Rudolph, The automorphism group of a shift of finite type, \emph{Trans. Amer. Math. Soc.} \textbf{306} (1988), 71--114.

\bibitem{BKP2021} B.~Bu\v{c}a, K.~Klobas and T.~Prosen, Rule 54: exactly solvable model of nonequilibrium statistical mechanics, \emph{J. Stat. Mech.} (2021), 074001.

\bibitem{CHR1979} E.~M.~Coven, G.~A.~Hedlund and F.~Rhodes, The commuting block maps problem, \emph{Trans. Amer. Math. Soc.} \textbf{249} (1979), 113--138.

\bibitem{CFK2019} V.~Cyr, J.~Franks and B.~Kra, The spacetime of a shift endomorphism, \emph{Trans. Amer. Math. Soc.} \textbf{371} (2019), 461--488.

\bibitem{CK2015} V.~Cyr and B.~Kra, The automorphism group of a shift of linear growth: beyond transitivity, \emph{Forum Math. Sigma} \textbf{3} (2015), e5.

\bibitem{PaperA} M.~Dayah, Reversible rule 110 has no hidden symmetries: six families of second-order cellular automata settled by edge moves, in preparation.

\bibitem{PaperB} M.~Dayah, Fronts at the speed of light: rigidity of permutive second-order cellular automata, in preparation.

\bibitem{PaperC} M.~Dayah, Every finite group from a wall: large automorphism groups of second-order cellular automata, in preparation.

\bibitem{PaperE} M.~Dayah, The cellular automata that commute with Rule 54, and rigidity of the second-order automaton 94R, in preparation.

\bibitem{DDMP2016} S.~Donoso, F.~Durand, A.~Maass and S.~Petite, On automorphism groups of low complexity subshifts, \emph{Ergodic Theory Dynam. Systems} \textbf{36} (2016), 64--95.

\bibitem{MiniSat} N.~E\'en and N.~S\"orensson, An extensible SAT-solver, in \emph{Theory and Applications of Satisfiability Testing (SAT 2003)}, Lecture Notes in Comput. Sci. 2919, Springer, 2004, 502--518.

\bibitem{HT1991} T.~Hattori and S.~Takesue, Additive conserved quantities in discrete-time lattice dynamical systems, \emph{Phys. D} \textbf{49} (1991), 295--322.

\bibitem{Hedlund1969} G.~A.~Hedlund, Endomorphisms and automorphisms of the shift dynamical system, \emph{Math. Systems Theory} \textbf{3} (1969), 320--375.

\bibitem{Hochman2010} M.~Hochman, On the automorphism groups of multidimensional shifts of finite type, \emph{Ergodic Theory Dynam. Systems} \textbf{30} (2010), 809--840.

\bibitem{PySAT} A.~Ignatiev, A.~Morgado and J.~Marques-Silva, PySAT: a Python toolkit for prototyping with SAT oracles, in \emph{Theory and Applications of Satisfiability Testing (SAT 2018)}, Lecture Notes in Comput. Sci. 10929, Springer, 2018, 428--437.

\bibitem{KimRoush1990} K.~H.~Kim and F.~W.~Roush, On the automorphism groups of subshifts, \emph{Pure Math. Appl. Ser. B} \textbf{1} (1990), no.~4, 203--230.

\bibitem{KS1989} B.~Kitchens and K.~Schmidt, Automorphisms of compact groups, \emph{Ergodic Theory Dynam. Systems} \textbf{9} (1989), 691--735.

\bibitem{KS1992} B.~Kitchens and K.~Schmidt, Markov subgroups of $(\Z/2\Z)^{\Z^2}$, in \emph{Symbolic Dynamics and its Applications}, Contemp. Math. 135, Amer. Math. Soc., 1992, 265--283.

\bibitem{KS2000} B.~Kitchens and K.~Schmidt, Isomorphism rigidity of irreducible algebraic $\mathbb{Z}^d$-actions, \emph{Invent. Math.} \textbf{142} (2000), 559--577.

\bibitem{Kopra2022} J.~Kopra, Glider automata on all transitive sofic shifts, \emph{Ergodic Theory Dynam. Systems} \textbf{42} (2022), 3716--3744.

\bibitem{Kurka1997} P.~K\r{u}rka, Languages, equicontinuity and attractors in cellular automata, \emph{Ergodic Theory Dynam. Systems} \textbf{17} (1997), 417--433.

\bibitem{Ledrappier1978} F.~Ledrappier, Un champ markovien peut \^etre d'entropie nulle et m\'elangeant, \emph{C. R. Acad. Sci. Paris S\'er. A-B} \textbf{287} (1978), A561--A563.

\bibitem{Margolus1984} N.~Margolus, Physics-like models of computation, \emph{Phys. D} \textbf{10} (1984), 81--95.

\bibitem{MACC2012} G.~J.~Mart\'inez, A.~Adamatzky, F.~Chen and L.~Chua, On soliton collisions between localizations in complex elementary cellular automata: rules 54 and 110 and beyond, \emph{Complex Systems} \textbf{21} (2012), 117--142; arXiv:1301.6258.

\bibitem{MooreBoykett1997} C.~Moore and T.~Boykett, Commuting cellular automata, \emph{Complex Systems} \textbf{11} (1997), 55--64.

\bibitem{Ryan1972} J.~P.~Ryan, The shift and commutativity, \emph{Math. Systems Theory} \textbf{6} (1972), 82--85.

\bibitem{Salo2019Ryan} V.~Salo, Transitive action on finite points of a full shift and a finitary Ryan's theorem, \emph{Ergodic Theory Dynam. Systems} \textbf{39} (2019), 1637--1667.

\bibitem{Salo2020} V.~Salo, Conjugacy of reversible cellular automata and one-head machines, preprint, arXiv:2011.07827 (2020).

\bibitem{SaloTorma2013} V.~Salo and I.~T\"orm\"a, Commutators of bipermutive and affine cellular automata, in \emph{Cellular Automata and Discrete Complex Systems (AUTOMATA 2013)}, Lecture Notes in Comput. Sci. 8155, Springer, 2013, 155--170.

\bibitem{Takesue1987} S.~Takesue, Reversible cellular automata and statistical mechanics, \emph{Phys. Rev. Lett.} \textbf{59} (1987), 2499--2502.

\bibitem{Takesue1989} S.~Takesue, Ergodic properties and thermodynamic behavior of elementary reversible cellular automata. I. Basic properties, \emph{J. Stat. Phys.} \textbf{56} (1989), 371--402.

\bibitem{Takesue1995} S.~Takesue, Staggered invariants in cellular automata, \emph{Complex Systems} \textbf{9} (1995), 149--168.

\bibitem{Takesue1997} S.~Takesue, Seru \=otomaton no hozonry\=o [Conserved quantities of cellular automata] (in Japanese), \emph{RIMS K\^oky\^uroku} \textbf{1020} (1997), 103--126.

\bibitem{TM1990} T.~Toffoli and N.~Margolus, Invertible cellular automata: a review, \emph{Phys. D} \textbf{45} (1990), 229--253.

\bibitem{Voorhees1993} B.~Voorhees, Commutation of cellular automata rules, \emph{Complex Systems} \textbf{7} (1993), 309--325.

\bibitem{Ward1994} T.~Ward, Automorphisms of $\mathbb{Z}^d$-subshifts of finite type, \emph{Indag. Math. (N.S.)} \textbf{5} (1994), 495--504.

\bibitem{Wolfram2002} S.~Wolfram, \emph{A New Kind of Science}, Wolfram Media, Champaign, IL, 2002.

\end{thebibliography}


\end{document}
